\begingroup
\tcbset{breakable/.default=false,after={\par\nobreak\smallskip}}
\RenewDocumentEnvironment{xesolution}{+b}{%
  \par\nobreak\smallskip\nobreak\begingroup\small
  \clubpenalties 2 10000 10000\displaywidowpenalties 2 10000 10000\noindent
  {\bfseries\color{DeepBlue}解答}\quad #1\par\endgroup\smallskip
}{}
\chapter{电磁场、电磁波与超导体}

{
\graphicspath{{/Users/dongxingxi/Documents/物理竞赛一轮资料/sample/xes-original/forty-first-section/}{/Users/dongxingxi/Documents/物理竞赛一轮资料/sample/xes-original/forty-first-section/assets/}}
\section{位移电流与麦克斯韦方程组}

% Diagrams for Lesson 10: electromagnetic fields and superconductors.
\newcommand{\maxwellOnnesContext}{%
  \historyfigure{maxwell-onnes-built-in-v1.png}%
  {左：麦克斯韦与电磁场理论；右：卡末林·昂内斯与低温超导实验。人物图用于建立历史情境，公式与物理结构见后文示意图。\label{ch26:ref:review-history-context}}}

\newcommand{\radialCurrentSphereDiagram}{%
\begin{center}
\begin{tikzpicture}[scale=.88,>=Latex]
 \draw[thick] (0,0) circle (1.7);
 \fill (0,0) circle (2.5pt);
 \foreach \a in {0,45,...,315}{\draw[vector] (\a:.8)--(\a:2.15);}
 \foreach \a in {22.5,67.5,...,337.5}{\draw[->,dashed,TextGray] (\a:.85)--(\a:1.95);}
 \node at (0,-.25) {$Q(t)$};
 \node[right] at (2.2,0) {$\vect E$};
 \node[above] at (.95,2) {$\vect j$};
 \begin{scope}[shift={(.9,.75)},rotate=25]
  \draw[DeepBlue] (0,0) ellipse (.42 and .22);
  \draw[->,DeepBlue] (.42,0) arc[start angle=0,end angle=95,x radius=.42,y radius=.22];
 \end{scope}
 \node[above] at (.8,1.08) {$\Gamma$};
 \node[below] at (0,-2.25) {球对称分布};
 \fill (5,0) circle (2pt) node[above] {$P$};
 \draw[vector] (5,-.1)--(6.7,-.1) node[right] {$\vect j$};
 \draw[->,TextGray,thick] (5,-.7)--(3.3,-.7) ;
 \node at (4.2,-1.2) {$\partial\vect D/\partial t$};
 \node[align=center] at (5,1.3) {同一点 $P$ 的两种源\\等大、反向};
 \node at (5,-2) {$\vect j+\partial\vect D/\partial t=0$};
\end{tikzpicture}
\figcap{径向泄漏电荷。右侧单独展示局部源的抵消。}\label{ch26:fig:s41-10-2}
\end{center}}

\newcommand{\displacementCurrentDiagram}{%
\begin{center}
\begin{tikzpicture}[scale=.9,>=Latex]
 % Two meridian views make the identical boundary and the deformed surface explicit.
 \foreach \xx/\tag in {0/S_1,7/S_2}{
  \begin{scope}[shift={(\xx,0)}]
   \draw[very thick] (-2.7,0)--(-.65,0) (.65,0)--(2.25,0);
   \draw[line width=2pt] (-.65,-1.1)--(-.65,1.1) (.65,-1.1)--(.65,1.1);
   \draw[->,DeepBlue,thick] (-2.6,.4)--(-1.9,.4) node[midway,above] {$I$};
   \fill[DeepBlue] (-1.7,.65) circle (2.2pt) (-1.7,-.65) circle (2.2pt);
   \node[left,DeepBlue] at (-1.4,.95) {$C$};
   \node[below] at (-.65,-1.8) {$+Q$};\node[below] at (.65,-1.8) {$-Q$};
  \end{scope}
 }
 \draw[DeepBlue,thick] (-1.7,-.65)--(-1.7,.65);
 \node[DeepBlue] at (-.6,-2.8) {$S_1$：平圆盘截线};
 \begin{scope}[shift={(7,0)}]
  \draw[DeepBlue,thick] (-1.7,.65)
   .. controls (-1.6,1.45) and (-1.2,1.7) .. (-.8,1.7)
   .. controls (-.25,1.7) and (0,1.45) .. (0,0);
  \draw[DeepBlue,thick] (-1.7,-.65)
   .. controls (-1.6,-1.45) and (-1.2,-1.7) .. (-.8,-1.7)
   .. controls (-.25,-1.7) and (0,-1.45) .. (0,0);
  \draw[->,TextGray] (-.45,.25)--(.45,.25);
  \node[right] at (1.1,.7) {$\partial\vect D/\partial t$};
  \node[DeepBlue] at (-.5,-2.8) {$S_2$：绕过左板后封口于间隙};
 \end{scope}
\end{tikzpicture}
\par\smallskip{\small\color{TextGray}同一轴截面的两种跨越面。蓝色剖面绕水平轴旋转成曲面，两个蓝点旋转成同一圆形边界 $C$；$S_2$ 不穿过极板或导线。}
\figcap{充电电容器的两个同边界跨越面。}\label{ch26:ref:review-surface-choice}
\end{center}}

\newcommand{\continuityDiagram}{%
\begin{center}
\begin{tikzpicture}[>=Latex]
 \node[draw=RuleGray,align=center,text width=8cm,inner sep=7pt] (a) at (0,1.7)
 {$\nabla\times\vect B=\mu_0\vect j$\\取散度只能得到 $\nabla\cdot\vect j=0$};
 \node[draw=DeepBlue,align=center,text width=8cm,inner sep=7pt] (b) at (0,-.3)
 {$\nabla\times\vect B=\mu_0(\vect j+\partial\vect D/\partial t)$\\
  配合 $\nabla\cdot\vect D=\rho$，得到 $\nabla\cdot\vect j=-\partial\rho/\partial t$};
 \draw[->,DeepBlue,thick] (a.south)--(b.north) node[midway,right] {补位移电流};
\end{tikzpicture}
\figcap{位移电流使环路定律与电荷连续性相容（非磁介质）。}\label{ch26:ref:review-continuity}
\end{center}}

\newcommand{\maxwellLightPanel}{%
\begin{center}
\begin{tikzpicture}[x=1cm,y=1cm,>=Latex]
  \node[draw=RuleGray,rounded corners=2pt,text width=8.8cm,
        minimum height=0.92cm,align=center,inner sep=5pt] (laws)
    at (0,1.75) {高斯定律、法拉第定律与安培--麦克斯韦定律};
  \node[draw=DeepBlue,rounded corners=2pt,text width=8.8cm,
        minimum height=1.08cm,align=center,inner sep=5pt] (wave)
    at (0,0) {在真空无源区联立，得到横向电磁波\\[-1pt]
      $\displaystyle c=\frac1{\sqrt{\mu_0\varepsilon_0}}$};
  \node[draw=DeepBlue,rounded corners=2pt,text width=8.8cm,
        minimum height=0.92cm,align=center,inner sep=5pt] (light)
    at (0,-1.82) {该速度在数值上等于真空光速 $c$：光是电磁波};
  \draw[->,DeepBlue,very thick] (laws.south)--(wave.north);
  \draw[->,DeepBlue,very thick] (wave.south)--(light.north);
\end{tikzpicture}
\figcap{由麦克斯韦方程到光的电磁理论。}\label{ch26:ref:review-maxwell-light}
\end{center}}

\newcommand{\planeWaveThreeDDiagram}{%
\begin{center}
\begin{tikzpicture}[x={(1.05cm,0cm)},y={(0cm,1cm)},z={(-0.46cm,-0.34cm)},>=Latex]
  \draw[axis] (0,0,0)--(8.25,0,0) node[right] {$x$（传播方向）};
  \draw[axis] (0,0,0)--(0,2.05,0) node[above] {$y$};
  \draw[axis] (0,0,0)--(0,0,2.4) node[below left] {$z$（朝向读者）};
  \draw[DeepBlue,very thick,domain=0:7.5,samples=120,smooth,variable=\t]
    plot ({\t},{1.18*sin(72*\t)},0);
  \draw[TextGray,very thick,domain=0:7.5,samples=120,smooth,variable=\t]
    plot ({\t},0,{1.05*sin(72*\t)});
  \foreach \t in {0.65,3.15,5.65}{
    \pgfmathsetmacro{\ey}{1.18*sin(72*\t)}
    \pgfmathsetmacro{\bz}{1.05*sin(72*\t)}
    \draw[->,DeepBlue,thick] (\t,0,0)--(\t,\ey,0);
    \draw[->,TextGray,thick] (\t,0,0)--(\t,0,\bz);
  }
  \node[DeepBlue,anchor=south west,inner sep=0.5pt] at (0.95,1.55,0.15) {$\vect E\parallel\uvec y$};
  \node[TextGray,anchor=south west,inner sep=0.5pt] at (.3,-1.45,0) {$\vect B\parallel\uvec z$};
  \draw[->,DeepBlue,line width=1.2pt] (5.55,1.72,0)--(7.4,1.72,0)
    node[midway,above] {$\vect S\parallel\vect k$};
  \node[DeepBlue,anchor=north east,inner sep=0.5pt] at (7.4,-1.72,0)
    {$\vect E\times\vect B\parallel +\uvec x$};
\end{tikzpicture}
\figcap{平面波的三维投影。电、磁场同相且与传播方向构成右手系。}\label{ch26:ref:review-wave-3d}
\end{center}}

\newcommand{\poyntingBalanceDiagram}{%
\begin{center}
\begin{tikzpicture}[>=Latex]
 \draw[draw=RuleGray,fill=SoftGray] (-2,-1.15) rectangle (2,1.15);
 \node at (0,.75) {体积 $V$};
 \draw[vector] (-1.2,0)--(.65,0) node[right] {$\vect E$};
 \draw[->,TextGray,thick] (-1.2,-.65)--(.65,-.65) node[right] {$\vect j$};
 \draw[vector] (2,.45)--(3.25,.45) node[right] {$\vect S\cdot\dd\vect A$};
 \draw[vector] (-2,0)--(-3,0);
 \node[align=center] at (0,-2)
 {场能减少率 $=$ 向外净能流 $+$ 对物质做功率\\[5pt]
 $\displaystyle-\frac{\dd}{\dd t}\int_Vu\,\dd V
 =\oint_{\partial V}\vect S\cdot\dd\vect A+\int_V\vect j\cdot\vect E\,\dd V$};
\end{tikzpicture}
\figcap{电磁场的能量流出与对物质做功。}\label{ch26:ref:review-energy-balance}
\end{center}}

\newcommand{\radiationPatternDiagram}{%
\begin{center}
\begin{tikzpicture}[scale=.83,>=Latex]
 \draw[axis] (0,-2.4)--(0,2.6) node[above] {$\vect a$};
 \draw[guide] (-2.9,0)--(2.9,0);
 \draw[DeepBlue,very thick] plot[domain=0:360,samples=240]
  ({2.6*sin(\x)^3},{2.6*cos(\x)*sin(\x)^2});
 \fill (0,0) circle (2.5pt);
 \node at (.35,-1.4) {$q$};
 \draw[->,TextGray] (0,0)--(40:2.7) node[above left] {$\uvec n$};
 \draw[->,DeepBlue,thick] (2.068,1.736)--(2.486,1.238);
 \node[above] at (3.4,2.6) {$\vect E_{\rm rad}\ (q>0)$};
 \draw (0,.7) arc[start angle=90,end angle=40,radius=.7];
 \node at (.3,1.3) {$\theta$};
 \node[align=left,anchor=west] at (3.4,.3)
 {蓝色径向长度 $\propto\dd P/\dd\Omega$\\[5pt]
  $\dd P/\dd\Omega\propto\sin^2\theta$\\[5pt]
  沿加速度轴为零，横向最强\\[5pt]
  曲线不是电荷轨迹或波前};
\end{tikzpicture}
\figcap{非相对论电荷的辐射功率角分布。}\label{ch26:ref:review-radiation-pattern}
\end{center}}

\newcommand{\radiationWavefrontDiagram}{%
\begin{center}
\begin{tikzpicture}[scale=.8,>=Latex]
 \draw[axis] (-4,0)--(4.1,0) node[right] {$r$};
 \draw[axis,dashed] (0,-3.6)--(0,3.7) node[above] {振动轴};
 \fill (0,0) circle (2.5pt);
 \node at (-.33,-.28) {$q$};
 \foreach \rad/\sgn in {.85/1,1.65/-1,2.45/1,3.25/-1}{
  \draw[dashed,TextGray] (0,0) circle (\rad);
  \foreach \a in {0,30,60,120,150,180,210,240,300,330}{
   \pgfmathsetmacro{\dx}{.32*\sgn*sin(\a)*cos(\a)/\rad}
   \pgfmathsetmacro{\dy}{-.32*\sgn*cos(\a)^2/\rad}
   \draw[->,DeepBlue,thick] ({\rad*cos(\a)-\dx},{\rad*sin(\a)-\dy})
    --({\rad*cos(\a)+\dx},{\rad*sin(\a)+\dy});
  }
 }
 \draw[<->,TextGray] (.85,-3.85)--(1.65,-3.85) node[midway,below] {$\lambda/2$};
 \draw[guide] (.85,0)--(.85,-3.85) (1.65,0)--(1.65,-3.85);
 \draw[vector] (3.3,.4)--(4.15,.4);
 \node[above] at (3.95,.5) {$\vect S$};
 \node[align=left,anchor=west] at (4.8,.2)
 {$q>0$；同一时刻，相邻圆相差半波长\\[5pt]
  箭头为横向电场，不是磁感线\\[5pt]
  大小同时体现 $\sin\theta$ 与 $1/r$\\[5pt]
  沿振动轴的电场箭头为零};
\end{tikzpicture}
\figcap{振荡电荷的远区辐射。相隔半波长的场反向，相隔一波长的场同向。}\label{ch26:fig:s41-10-7}
\end{center}}

\newcommand{\thomsonScatteringDiagram}{%
\begin{center}
\begin{tikzpicture}[scale=.85,>=Latex]
 \draw[axis] (-4,0)--(4,0) node[right] {$x$};
 \draw[DeepBlue,thick] plot[domain=-3.9:-1.3,samples=70] (\x,{.45*sin(180*\x)});
 \draw[vector] (-3.7,1.2)--(-2.3,1.2) node[midway,above] {入射波};
 \fill (0,0) circle (2.5pt);
 \node[anchor=east] at (-.45,-.28) {$e^-$};
 \draw[vector] (0,0)--(0,1.3) node[above] {$\vect E$};
 \draw[->,TextGray,thick] (0,0)--(0,-1.3) node[below] {$\vect a$};
 \foreach \a in {-50,0,50,130,180,230}{
  \draw[->,TextGray] (\a:.7)--(\a:2);
 }
 \node[align=left,anchor=west] at (4.6,.2)
 {选取 $E_z>0$ 的瞬间\\[4pt]
  $a_z=-eE_z/m<0$\\[4pt]
  横向辐射强，振动轴上无辐射\\[4pt]
  $\sigma_T=(8\pi/3)r_e^2$};
\end{tikzpicture}
\figcap{电子受场驱动的瞬时方向与再辐射。}\label{ch26:ref:review-thomson}
\end{center}}

\newcommand{\superconductivityTimeline}{%
\begin{center}
\begin{tikzpicture}[x=0.86cm,y=1cm,>=Latex]
  \draw[axis] (-0.2,0)--(13.3,0) node[right] {时间};
  \foreach \x/\year/\desc in {
    0.7/1911/{Hg，$4.2\,$K},
    3.55/1933/{迈斯纳效应},
    6.4/1987/{铜氧化物，$T_c>77\,$K},
    9.25/2001/{MgB$_2$，$39\,$K},
    12.1/2015/{高压硫氢化物}}
  {
    \fill[DeepBlue] (\x,0) circle (2.6pt);
    \draw[TextGray] (\x,0.12)--(\x,0.7);
    \node[align=center,anchor=south,font=\scriptsize,text width=2.25cm] at (\x,0.72) {\year\\\desc};
  }
  \node[TextGray,align=left,anchor=north west,font=\small,text width=13.5cm] at (0.65,-0.38)
    {这是一条概念史线索，不表示“临界温度越高，材料越容易应用”；临界磁场、电流密度与制备条件同样关键。};
\end{tikzpicture}
\figcap{超导研究的历史节点。}\label{ch26:ref:review-superconductivity-timeline}
\end{center}}

\newcommand{\meissnerComparisonDiagram}{%
\begin{center}
\begin{tikzpicture}[scale=.82,>=Latex]
 \foreach \cx/\title in {-3.6/{理想导体：冻结初始场},3.6/{迈斯纳态：排斥内部场}}{
  \draw[thick,fill=PaleBlue] (\cx,0) circle (1.15);
  \node[above] at (\cx,2.8) {\title};
 }
 \foreach \x in {-1.5,-1,-.5,0,.5,1,1.5}{
  \draw[->,TextGray] ({-3.6+\x},-2.4)--({-3.6+\x},2.4);
 }
 % Analytic field lines outside a perfectly diamagnetic cylinder:
 % psi/B0 = x*(1-a^2/(x^2+y^2)); solve by bisection before exporting.
 \draw[->,DeepBlue,thick] plot coordinates {(3.444431,-2.400000) (3.442849,-2.360000) (3.441152,-2.320000) (3.439326,-2.280000) (3.437357,-2.240000) (3.435228,-2.200000) (3.432922,-2.160000) (3.430416,-2.120000) (3.427686,-2.080000) (3.424701,-2.040000) (3.421426,-2.000000) (3.417821,-1.960000) (3.413835,-1.920000) (3.409411,-1.880000) (3.404476,-1.840000) (3.398945,-1.800000) (3.392710,-1.760000) (3.385640,-1.720000) (3.377572,-1.680000) (3.368301,-1.640000) (3.357572,-1.600000) (3.345062,-1.560000) (3.330368,-1.520000) (3.312995,-1.480000) (3.292355,-1.440000) (3.267795,-1.400000) (3.238686,-1.360000) (3.204596,-1.320000) (3.165532,-1.280000) (3.122131,-1.240000) (3.075631,-1.200000) (3.027572,-1.160000) (2.979413,-1.120000) (2.932282,-1.080000) (2.886918,-1.040000) (2.843741,-1.000000) (2.802945,-0.960000) (2.764579,-0.920000) (2.728616,-0.880000) (2.694980,-0.840000) (2.663579,-0.800000) (2.634310,-0.760000) (2.607074,-0.720000) (2.581775,-0.680000) (2.558325,-0.640000) (2.536643,-0.600000) (2.516654,-0.560000) (2.498293,-0.520000) (2.481502,-0.480000) (2.466226,-0.440000) (2.452420,-0.400000) (2.440042,-0.360000) (2.429057,-0.320000) (2.419434,-0.280000) (2.411145,-0.240000) (2.404168,-0.200000) (2.398485,-0.160000) (2.394080,-0.120000) (2.390942,-0.080000) (2.389062,-0.040000) (2.388436,0.000000) (2.389062,0.040000) (2.390942,0.080000) (2.394080,0.120000) (2.398485,0.160000) (2.404168,0.200000) (2.411145,0.240000) (2.419434,0.280000) (2.429057,0.320000) (2.440042,0.360000) (2.452420,0.400000) (2.466226,0.440000) (2.481502,0.480000) (2.498293,0.520000) (2.516654,0.560000) (2.536643,0.600000) (2.558325,0.640000) (2.581775,0.680000) (2.607074,0.720000) (2.634310,0.760000) (2.663579,0.800000) (2.694980,0.840000) (2.728616,0.880000) (2.764579,0.920000) (2.802945,0.960000) (2.843741,1.000000) (2.886918,1.040000) (2.932282,1.080000) (2.979413,1.120000) (3.027572,1.160000) (3.075631,1.200000) (3.122131,1.240000) (3.165532,1.280000) (3.204596,1.320000) (3.238686,1.360000) (3.267795,1.400000) (3.292355,1.440000) (3.312995,1.480000) (3.330368,1.520000) (3.345062,1.560000) (3.357572,1.600000) (3.368301,1.640000) (3.377572,1.680000) (3.385640,1.720000) (3.392710,1.760000) (3.398945,1.800000) (3.404476,1.840000) (3.409411,1.880000) (3.413835,1.920000) (3.417821,1.960000) (3.421426,2.000000) (3.424701,2.040000) (3.427686,2.080000) (3.430416,2.120000) (3.432922,2.160000) (3.435228,2.200000) (3.437357,2.240000) (3.439326,2.280000) (3.441152,2.320000) (3.442849,2.360000) (3.444431,2.400000)};
\draw[->,DeepBlue,thick] plot coordinates {(3.150238,-2.400000) (3.146057,-2.360000) (3.141599,-2.320000) (3.136839,-2.280000) (3.131749,-2.240000) (3.126299,-2.200000) (3.120454,-2.160000) (3.114177,-2.120000) (3.107425,-2.080000) (3.100151,-2.040000) (3.092305,-2.000000) (3.083829,-1.960000) (3.074661,-1.920000) (3.064732,-1.880000) (3.053970,-1.840000) (3.042296,-1.800000) (3.029628,-1.760000) (3.015881,-1.720000) (3.000973,-1.680000) (2.984824,-1.640000) (2.967367,-1.600000) (2.948549,-1.560000) (2.928341,-1.520000) (2.906748,-1.480000) (2.883810,-1.440000) (2.859611,-1.400000) (2.834279,-1.360000) (2.807984,-1.320000) (2.780925,-1.280000) (2.753328,-1.240000) (2.725425,-1.200000) (2.697447,-1.160000) (2.669610,-1.120000) (2.642113,-1.080000) (2.615124,-1.040000) (2.588791,-1.000000) (2.563231,-0.960000) (2.538540,-0.920000) (2.514790,-0.880000) (2.492037,-0.840000) (2.470321,-0.800000) (2.449671,-0.760000) (2.430105,-0.720000) (2.411633,-0.680000) (2.394261,-0.640000) (2.377988,-0.600000) (2.362812,-0.560000) (2.348729,-0.520000) (2.335730,-0.480000) (2.323809,-0.440000) (2.312959,-0.400000) (2.303170,-0.360000) (2.294435,-0.320000) (2.286747,-0.280000) (2.280099,-0.240000) (2.274484,-0.200000) (2.269898,-0.160000) (2.266335,-0.120000) (2.263793,-0.080000) (2.262269,-0.040000) (2.261761,0.000000) (2.262269,0.040000) (2.263793,0.080000) (2.266335,0.120000) (2.269898,0.160000) (2.274484,0.200000) (2.280099,0.240000) (2.286747,0.280000) (2.294435,0.320000) (2.303170,0.360000) (2.312959,0.400000) (2.323809,0.440000) (2.335730,0.480000) (2.348729,0.520000) (2.362812,0.560000) (2.377988,0.600000) (2.394261,0.640000) (2.411633,0.680000) (2.430105,0.720000) (2.449671,0.760000) (2.470321,0.800000) (2.492037,0.840000) (2.514790,0.880000) (2.538540,0.920000) (2.563231,0.960000) (2.588791,1.000000) (2.615124,1.040000) (2.642113,1.080000) (2.669610,1.120000) (2.697447,1.160000) (2.725425,1.200000) (2.753328,1.240000) (2.780925,1.280000) (2.807984,1.320000) (2.834279,1.360000) (2.859611,1.400000) (2.883810,1.440000) (2.906748,1.480000) (2.928341,1.520000) (2.948549,1.560000) (2.967367,1.600000) (2.984824,1.640000) (3.000973,1.680000) (3.015881,1.720000) (3.029628,1.760000) (3.042296,1.800000) (3.053970,1.840000) (3.064732,1.880000) (3.074661,1.920000) (3.083829,1.960000) (3.092305,2.000000) (3.100151,2.040000) (3.107425,2.080000) (3.114177,2.120000) (3.120454,2.160000) (3.126299,2.200000) (3.131749,2.240000) (3.136839,2.280000) (3.141599,2.320000) (3.146057,2.360000) (3.150238,2.400000)};
\draw[->,DeepBlue,thick] plot coordinates {(2.722243,-2.400000) (2.715907,-2.360000) (2.709255,-2.320000) (2.702267,-2.280000) (2.694927,-2.240000) (2.687217,-2.200000) (2.679117,-2.160000) (2.670608,-2.120000) (2.661672,-2.080000) (2.652291,-2.040000) (2.642447,-2.000000) (2.632124,-1.960000) (2.621306,-1.920000) (2.609981,-1.880000) (2.598140,-1.840000) (2.585776,-1.800000) (2.572887,-1.760000) (2.559475,-1.720000) (2.545547,-1.680000) (2.531118,-1.640000) (2.516207,-1.600000) (2.500841,-1.560000) (2.485053,-1.520000) (2.468883,-1.480000) (2.452376,-1.440000) (2.435584,-1.400000) (2.418564,-1.360000) (2.401374,-1.320000) (2.384079,-1.280000) (2.366742,-1.240000) (2.349427,-1.200000) (2.332200,-1.160000) (2.315123,-1.120000) (2.298255,-1.080000) (2.281653,-1.040000) (2.265373,-1.000000) (2.249462,-0.960000) (2.233967,-0.920000) (2.218929,-0.880000) (2.204386,-0.840000) (2.190371,-0.800000) (2.176913,-0.760000) (2.164040,-0.720000) (2.151774,-0.680000) (2.140134,-0.640000) (2.129139,-0.600000) (2.118804,-0.560000) (2.109141,-0.520000) (2.100160,-0.480000) (2.091872,-0.440000) (2.084285,-0.400000) (2.077404,-0.360000) (2.071235,-0.320000) (2.065783,-0.280000) (2.061051,-0.240000) (2.057043,-0.200000) (2.053760,-0.160000) (2.051205,-0.120000) (2.049380,-0.080000) (2.048284,-0.040000) (2.047918,0.000000) (2.048284,0.040000) (2.049380,0.080000) (2.051205,0.120000) (2.053760,0.160000) (2.057043,0.200000) (2.061051,0.240000) (2.065783,0.280000) (2.071235,0.320000) (2.077404,0.360000) (2.084285,0.400000) (2.091872,0.440000) (2.100160,0.480000) (2.109141,0.520000) (2.118804,0.560000) (2.129139,0.600000) (2.140134,0.640000) (2.151774,0.680000) (2.164040,0.720000) (2.176913,0.760000) (2.190371,0.800000) (2.204386,0.840000) (2.218929,0.880000) (2.233967,0.920000) (2.249462,0.960000) (2.265373,1.000000) (2.281653,1.040000) (2.298255,1.080000) (2.315123,1.120000) (2.332200,1.160000) (2.349427,1.200000) (2.366742,1.240000) (2.384079,1.280000) (2.401374,1.320000) (2.418564,1.360000) (2.435584,1.400000) (2.452376,1.440000) (2.468883,1.480000) (2.485053,1.520000) (2.500841,1.560000) (2.516207,1.600000) (2.531118,1.640000) (2.545547,1.680000) (2.559475,1.720000) (2.572887,1.760000) (2.585776,1.800000) (2.598140,1.840000) (2.609981,1.880000) (2.621306,1.920000) (2.632124,1.960000) (2.642447,2.000000) (2.652291,2.040000) (2.661672,2.080000) (2.670608,2.120000) (2.679117,2.160000) (2.687217,2.200000) (2.694927,2.240000) (2.702267,2.280000) (2.709255,2.320000) (2.715907,2.360000) (2.722243,2.400000)};
\draw[->,DeepBlue,thick] plot coordinates {(2.266006,-2.400000) (2.259327,-2.360000) (2.252420,-2.320000) (2.245280,-2.280000) (2.237902,-2.240000) (2.230283,-2.200000) (2.222420,-2.160000) (2.214309,-2.120000) (2.205949,-2.080000) (2.197340,-2.040000) (2.188482,-2.000000) (2.179378,-1.960000) (2.170029,-1.920000) (2.160441,-1.880000) (2.150620,-1.840000) (2.140574,-1.800000) (2.130313,-1.760000) (2.119847,-1.720000) (2.109191,-1.680000) (2.098360,-1.640000) (2.087369,-1.600000) (2.076239,-1.560000) (2.064988,-1.520000) (2.053639,-1.480000) (2.042214,-1.440000) (2.030738,-1.400000) (2.019236,-1.360000) (2.007734,-1.320000) (1.996258,-1.280000) (1.984836,-1.240000) (1.973493,-1.200000) (1.962259,-1.160000) (1.951158,-1.120000) (1.940218,-1.080000) (1.929463,-1.040000) (1.918920,-1.000000) (1.908612,-0.960000) (1.898563,-0.920000) (1.888793,-0.880000) (1.879325,-0.840000) (1.870179,-0.800000) (1.861372,-0.760000) (1.852922,-0.720000) (1.844845,-0.680000) (1.837157,-0.640000) (1.829872,-0.600000) (1.823002,-0.560000) (1.816558,-0.520000) (1.810552,-0.480000) (1.804994,-0.440000) (1.799891,-0.400000) (1.795252,-0.360000) (1.791084,-0.320000) (1.787392,-0.280000) (1.784183,-0.240000) (1.781459,-0.200000) (1.779226,-0.160000) (1.777486,-0.120000) (1.776242,-0.080000) (1.775494,-0.040000) (1.775245,0.000000) (1.775494,0.040000) (1.776242,0.080000) (1.777486,0.120000) (1.779226,0.160000) (1.781459,0.200000) (1.784183,0.240000) (1.787392,0.280000) (1.791084,0.320000) (1.795252,0.360000) (1.799891,0.400000) (1.804994,0.440000) (1.810552,0.480000) (1.816558,0.520000) (1.823002,0.560000) (1.829872,0.600000) (1.837157,0.640000) (1.844845,0.680000) (1.852922,0.720000) (1.861372,0.760000) (1.870179,0.800000) (1.879325,0.840000) (1.888793,0.880000) (1.898563,0.920000) (1.908612,0.960000) (1.918920,1.000000) (1.929463,1.040000) (1.940218,1.080000) (1.951158,1.120000) (1.962259,1.160000) (1.973493,1.200000) (1.984836,1.240000) (1.996258,1.280000) (2.007734,1.320000) (2.019236,1.360000) (2.030738,1.400000) (2.042214,1.440000) (2.053639,1.480000) (2.064988,1.520000) (2.076239,1.560000) (2.087369,1.600000) (2.098360,1.640000) (2.109191,1.680000) (2.119847,1.720000) (2.130313,1.760000) (2.140574,1.800000) (2.150620,1.840000) (2.160441,1.880000) (2.170029,1.920000) (2.179378,1.960000) (2.188482,2.000000) (2.197340,2.040000) (2.205949,2.080000) (2.214309,2.120000) (2.222420,2.160000) (2.230283,2.200000) (2.237902,2.240000) (2.245280,2.280000) (2.252420,2.320000) (2.259327,2.360000) (2.266006,2.400000)};
\draw[->,DeepBlue,thick] plot coordinates {(1.784889,-2.400000) (1.778969,-2.360000) (1.772920,-2.320000) (1.766741,-2.280000) (1.760433,-2.240000) (1.753999,-2.200000) (1.747439,-2.160000) (1.740757,-2.120000) (1.733955,-2.080000) (1.727037,-2.040000) (1.720007,-2.000000) (1.712870,-1.960000) (1.705630,-1.920000) (1.698294,-1.880000) (1.690869,-1.840000) (1.683362,-1.800000) (1.675780,-1.760000) (1.668132,-1.720000) (1.660427,-1.680000) (1.652674,-1.640000) (1.644884,-1.600000) (1.637068,-1.560000) (1.629237,-1.520000) (1.621402,-1.480000) (1.613576,-1.440000) (1.605771,-1.400000) (1.597999,-1.360000) (1.590273,-1.320000) (1.582608,-1.280000) (1.575015,-1.240000) (1.567509,-1.200000) (1.560102,-1.160000) (1.552807,-1.120000) (1.545639,-1.080000) (1.538609,-1.040000) (1.531731,-1.000000) (1.525017,-0.960000) (1.518480,-0.920000) (1.512130,-0.880000) (1.505981,-0.840000) (1.500041,-0.800000) (1.494324,-0.760000) (1.488837,-0.720000) (1.483592,-0.680000) (1.478598,-0.640000) (1.473862,-0.600000) (1.469394,-0.560000) (1.465201,-0.520000) (1.461290,-0.480000) (1.457668,-0.440000) (1.454341,-0.400000) (1.451313,-0.360000) (1.448592,-0.320000) (1.446180,-0.280000) (1.444081,-0.240000) (1.442300,-0.200000) (1.440839,-0.160000) (1.439700,-0.120000) (1.438885,-0.080000) (1.438395,-0.040000) (1.438232,0.000000) (1.438395,0.040000) (1.438885,0.080000) (1.439700,0.120000) (1.440839,0.160000) (1.442300,0.200000) (1.444081,0.240000) (1.446180,0.280000) (1.448592,0.320000) (1.451313,0.360000) (1.454341,0.400000) (1.457668,0.440000) (1.461290,0.480000) (1.465201,0.520000) (1.469394,0.560000) (1.473862,0.600000) (1.478598,0.640000) (1.483592,0.680000) (1.488837,0.720000) (1.494324,0.760000) (1.500041,0.800000) (1.505981,0.840000) (1.512130,0.880000) (1.518480,0.920000) (1.525017,0.960000) (1.531731,1.000000) (1.538609,1.040000) (1.545639,1.080000) (1.552807,1.120000) (1.560102,1.160000) (1.567509,1.200000) (1.575015,1.240000) (1.582608,1.280000) (1.590273,1.320000) (1.597999,1.360000) (1.605771,1.400000) (1.613576,1.440000) (1.621402,1.480000) (1.629237,1.520000) (1.637068,1.560000) (1.644884,1.600000) (1.652674,1.640000) (1.660427,1.680000) (1.668132,1.720000) (1.675780,1.760000) (1.683362,1.800000) (1.690869,1.840000) (1.698294,1.880000) (1.705630,1.920000) (1.712870,1.960000) (1.720007,2.000000) (1.727037,2.040000) (1.733955,2.080000) (1.740757,2.120000) (1.747439,2.160000) (1.753999,2.200000) (1.760433,2.240000) (1.766741,2.280000) (1.772920,2.320000) (1.778969,2.360000) (1.784889,2.400000)};
\draw[->,DeepBlue,thick] plot coordinates {(3.755569,-2.400000) (3.757151,-2.360000) (3.758848,-2.320000) (3.760674,-2.280000) (3.762643,-2.240000) (3.764772,-2.200000) (3.767078,-2.160000) (3.769584,-2.120000) (3.772314,-2.080000) (3.775299,-2.040000) (3.778574,-2.000000) (3.782179,-1.960000) (3.786165,-1.920000) (3.790589,-1.880000) (3.795524,-1.840000) (3.801055,-1.800000) (3.807290,-1.760000) (3.814360,-1.720000) (3.822428,-1.680000) (3.831699,-1.640000) (3.842428,-1.600000) (3.854938,-1.560000) (3.869632,-1.520000) (3.887005,-1.480000) (3.907645,-1.440000) (3.932205,-1.400000) (3.961314,-1.360000) (3.995404,-1.320000) (4.034468,-1.280000) (4.077869,-1.240000) (4.124369,-1.200000) (4.172428,-1.160000) (4.220587,-1.120000) (4.267718,-1.080000) (4.313082,-1.040000) (4.356259,-1.000000) (4.397055,-0.960000) (4.435421,-0.920000) (4.471384,-0.880000) (4.505020,-0.840000) (4.536421,-0.800000) (4.565690,-0.760000) (4.592926,-0.720000) (4.618225,-0.680000) (4.641675,-0.640000) (4.663357,-0.600000) (4.683346,-0.560000) (4.701707,-0.520000) (4.718498,-0.480000) (4.733774,-0.440000) (4.747580,-0.400000) (4.759958,-0.360000) (4.770943,-0.320000) (4.780566,-0.280000) (4.788855,-0.240000) (4.795832,-0.200000) (4.801515,-0.160000) (4.805920,-0.120000) (4.809058,-0.080000) (4.810938,-0.040000) (4.811564,0.000000) (4.810938,0.040000) (4.809058,0.080000) (4.805920,0.120000) (4.801515,0.160000) (4.795832,0.200000) (4.788855,0.240000) (4.780566,0.280000) (4.770943,0.320000) (4.759958,0.360000) (4.747580,0.400000) (4.733774,0.440000) (4.718498,0.480000) (4.701707,0.520000) (4.683346,0.560000) (4.663357,0.600000) (4.641675,0.640000) (4.618225,0.680000) (4.592926,0.720000) (4.565690,0.760000) (4.536421,0.800000) (4.505020,0.840000) (4.471384,0.880000) (4.435421,0.920000) (4.397055,0.960000) (4.356259,1.000000) (4.313082,1.040000) (4.267718,1.080000) (4.220587,1.120000) (4.172428,1.160000) (4.124369,1.200000) (4.077869,1.240000) (4.034468,1.280000) (3.995404,1.320000) (3.961314,1.360000) (3.932205,1.400000) (3.907645,1.440000) (3.887005,1.480000) (3.869632,1.520000) (3.854938,1.560000) (3.842428,1.600000) (3.831699,1.640000) (3.822428,1.680000) (3.814360,1.720000) (3.807290,1.760000) (3.801055,1.800000) (3.795524,1.840000) (3.790589,1.880000) (3.786165,1.920000) (3.782179,1.960000) (3.778574,2.000000) (3.775299,2.040000) (3.772314,2.080000) (3.769584,2.120000) (3.767078,2.160000) (3.764772,2.200000) (3.762643,2.240000) (3.760674,2.280000) (3.758848,2.320000) (3.757151,2.360000) (3.755569,2.400000)};
\draw[->,DeepBlue,thick] plot coordinates {(4.049762,-2.400000) (4.053943,-2.360000) (4.058401,-2.320000) (4.063161,-2.280000) (4.068251,-2.240000) (4.073701,-2.200000) (4.079546,-2.160000) (4.085823,-2.120000) (4.092575,-2.080000) (4.099849,-2.040000) (4.107695,-2.000000) (4.116171,-1.960000) (4.125339,-1.920000) (4.135268,-1.880000) (4.146030,-1.840000) (4.157704,-1.800000) (4.170372,-1.760000) (4.184119,-1.720000) (4.199027,-1.680000) (4.215176,-1.640000) (4.232633,-1.600000) (4.251451,-1.560000) (4.271659,-1.520000) (4.293252,-1.480000) (4.316190,-1.440000) (4.340389,-1.400000) (4.365721,-1.360000) (4.392016,-1.320000) (4.419075,-1.280000) (4.446672,-1.240000) (4.474575,-1.200000) (4.502553,-1.160000) (4.530390,-1.120000) (4.557887,-1.080000) (4.584876,-1.040000) (4.611209,-1.000000) (4.636769,-0.960000) (4.661460,-0.920000) (4.685210,-0.880000) (4.707963,-0.840000) (4.729679,-0.800000) (4.750329,-0.760000) (4.769895,-0.720000) (4.788367,-0.680000) (4.805739,-0.640000) (4.822012,-0.600000) (4.837188,-0.560000) (4.851271,-0.520000) (4.864270,-0.480000) (4.876191,-0.440000) (4.887041,-0.400000) (4.896830,-0.360000) (4.905565,-0.320000) (4.913253,-0.280000) (4.919901,-0.240000) (4.925516,-0.200000) (4.930102,-0.160000) (4.933665,-0.120000) (4.936207,-0.080000) (4.937731,-0.040000) (4.938239,0.000000) (4.937731,0.040000) (4.936207,0.080000) (4.933665,0.120000) (4.930102,0.160000) (4.925516,0.200000) (4.919901,0.240000) (4.913253,0.280000) (4.905565,0.320000) (4.896830,0.360000) (4.887041,0.400000) (4.876191,0.440000) (4.864270,0.480000) (4.851271,0.520000) (4.837188,0.560000) (4.822012,0.600000) (4.805739,0.640000) (4.788367,0.680000) (4.769895,0.720000) (4.750329,0.760000) (4.729679,0.800000) (4.707963,0.840000) (4.685210,0.880000) (4.661460,0.920000) (4.636769,0.960000) (4.611209,1.000000) (4.584876,1.040000) (4.557887,1.080000) (4.530390,1.120000) (4.502553,1.160000) (4.474575,1.200000) (4.446672,1.240000) (4.419075,1.280000) (4.392016,1.320000) (4.365721,1.360000) (4.340389,1.400000) (4.316190,1.440000) (4.293252,1.480000) (4.271659,1.520000) (4.251451,1.560000) (4.232633,1.600000) (4.215176,1.640000) (4.199027,1.680000) (4.184119,1.720000) (4.170372,1.760000) (4.157704,1.800000) (4.146030,1.840000) (4.135268,1.880000) (4.125339,1.920000) (4.116171,1.960000) (4.107695,2.000000) (4.099849,2.040000) (4.092575,2.080000) (4.085823,2.120000) (4.079546,2.160000) (4.073701,2.200000) (4.068251,2.240000) (4.063161,2.280000) (4.058401,2.320000) (4.053943,2.360000) (4.049762,2.400000)};
\draw[->,DeepBlue,thick] plot coordinates {(4.477757,-2.400000) (4.484093,-2.360000) (4.490745,-2.320000) (4.497733,-2.280000) (4.505073,-2.240000) (4.512783,-2.200000) (4.520883,-2.160000) (4.529392,-2.120000) (4.538328,-2.080000) (4.547709,-2.040000) (4.557553,-2.000000) (4.567876,-1.960000) (4.578694,-1.920000) (4.590019,-1.880000) (4.601860,-1.840000) (4.614224,-1.800000) (4.627113,-1.760000) (4.640525,-1.720000) (4.654453,-1.680000) (4.668882,-1.640000) (4.683793,-1.600000) (4.699159,-1.560000) (4.714947,-1.520000) (4.731117,-1.480000) (4.747624,-1.440000) (4.764416,-1.400000) (4.781436,-1.360000) (4.798626,-1.320000) (4.815921,-1.280000) (4.833258,-1.240000) (4.850573,-1.200000) (4.867800,-1.160000) (4.884877,-1.120000) (4.901745,-1.080000) (4.918347,-1.040000) (4.934627,-1.000000) (4.950538,-0.960000) (4.966033,-0.920000) (4.981071,-0.880000) (4.995614,-0.840000) (5.009629,-0.800000) (5.023087,-0.760000) (5.035960,-0.720000) (5.048226,-0.680000) (5.059866,-0.640000) (5.070861,-0.600000) (5.081196,-0.560000) (5.090859,-0.520000) (5.099840,-0.480000) (5.108128,-0.440000) (5.115715,-0.400000) (5.122596,-0.360000) (5.128765,-0.320000) (5.134217,-0.280000) (5.138949,-0.240000) (5.142957,-0.200000) (5.146240,-0.160000) (5.148795,-0.120000) (5.150620,-0.080000) (5.151716,-0.040000) (5.152082,0.000000) (5.151716,0.040000) (5.150620,0.080000) (5.148795,0.120000) (5.146240,0.160000) (5.142957,0.200000) (5.138949,0.240000) (5.134217,0.280000) (5.128765,0.320000) (5.122596,0.360000) (5.115715,0.400000) (5.108128,0.440000) (5.099840,0.480000) (5.090859,0.520000) (5.081196,0.560000) (5.070861,0.600000) (5.059866,0.640000) (5.048226,0.680000) (5.035960,0.720000) (5.023087,0.760000) (5.009629,0.800000) (4.995614,0.840000) (4.981071,0.880000) (4.966033,0.920000) (4.950538,0.960000) (4.934627,1.000000) (4.918347,1.040000) (4.901745,1.080000) (4.884877,1.120000) (4.867800,1.160000) (4.850573,1.200000) (4.833258,1.240000) (4.815921,1.280000) (4.798626,1.320000) (4.781436,1.360000) (4.764416,1.400000) (4.747624,1.440000) (4.731117,1.480000) (4.714947,1.520000) (4.699159,1.560000) (4.683793,1.600000) (4.668882,1.640000) (4.654453,1.680000) (4.640525,1.720000) (4.627113,1.760000) (4.614224,1.800000) (4.601860,1.840000) (4.590019,1.880000) (4.578694,1.920000) (4.567876,1.960000) (4.557553,2.000000) (4.547709,2.040000) (4.538328,2.080000) (4.529392,2.120000) (4.520883,2.160000) (4.512783,2.200000) (4.505073,2.240000) (4.497733,2.280000) (4.490745,2.320000) (4.484093,2.360000) (4.477757,2.400000)};
\draw[->,DeepBlue,thick] plot coordinates {(4.933994,-2.400000) (4.940673,-2.360000) (4.947580,-2.320000) (4.954720,-2.280000) (4.962098,-2.240000) (4.969717,-2.200000) (4.977580,-2.160000) (4.985691,-2.120000) (4.994051,-2.080000) (5.002660,-2.040000) (5.011518,-2.000000) (5.020622,-1.960000) (5.029971,-1.920000) (5.039559,-1.880000) (5.049380,-1.840000) (5.059426,-1.800000) (5.069687,-1.760000) (5.080153,-1.720000) (5.090809,-1.680000) (5.101640,-1.640000) (5.112631,-1.600000) (5.123761,-1.560000) (5.135012,-1.520000) (5.146361,-1.480000) (5.157786,-1.440000) (5.169262,-1.400000) (5.180764,-1.360000) (5.192266,-1.320000) (5.203742,-1.280000) (5.215164,-1.240000) (5.226507,-1.200000) (5.237741,-1.160000) (5.248842,-1.120000) (5.259782,-1.080000) (5.270537,-1.040000) (5.281080,-1.000000) (5.291388,-0.960000) (5.301437,-0.920000) (5.311207,-0.880000) (5.320675,-0.840000) (5.329821,-0.800000) (5.338628,-0.760000) (5.347078,-0.720000) (5.355155,-0.680000) (5.362843,-0.640000) (5.370128,-0.600000) (5.376998,-0.560000) (5.383442,-0.520000) (5.389448,-0.480000) (5.395006,-0.440000) (5.400109,-0.400000) (5.404748,-0.360000) (5.408916,-0.320000) (5.412608,-0.280000) (5.415817,-0.240000) (5.418541,-0.200000) (5.420774,-0.160000) (5.422514,-0.120000) (5.423758,-0.080000) (5.424506,-0.040000) (5.424755,0.000000) (5.424506,0.040000) (5.423758,0.080000) (5.422514,0.120000) (5.420774,0.160000) (5.418541,0.200000) (5.415817,0.240000) (5.412608,0.280000) (5.408916,0.320000) (5.404748,0.360000) (5.400109,0.400000) (5.395006,0.440000) (5.389448,0.480000) (5.383442,0.520000) (5.376998,0.560000) (5.370128,0.600000) (5.362843,0.640000) (5.355155,0.680000) (5.347078,0.720000) (5.338628,0.760000) (5.329821,0.800000) (5.320675,0.840000) (5.311207,0.880000) (5.301437,0.920000) (5.291388,0.960000) (5.281080,1.000000) (5.270537,1.040000) (5.259782,1.080000) (5.248842,1.120000) (5.237741,1.160000) (5.226507,1.200000) (5.215164,1.240000) (5.203742,1.280000) (5.192266,1.320000) (5.180764,1.360000) (5.169262,1.400000) (5.157786,1.440000) (5.146361,1.480000) (5.135012,1.520000) (5.123761,1.560000) (5.112631,1.600000) (5.101640,1.640000) (5.090809,1.680000) (5.080153,1.720000) (5.069687,1.760000) (5.059426,1.800000) (5.049380,1.840000) (5.039559,1.880000) (5.029971,1.920000) (5.020622,1.960000) (5.011518,2.000000) (5.002660,2.040000) (4.994051,2.080000) (4.985691,2.120000) (4.977580,2.160000) (4.969717,2.200000) (4.962098,2.240000) (4.954720,2.280000) (4.947580,2.320000) (4.940673,2.360000) (4.933994,2.400000)};
\draw[->,DeepBlue,thick] plot coordinates {(5.415111,-2.400000) (5.421031,-2.360000) (5.427080,-2.320000) (5.433259,-2.280000) (5.439567,-2.240000) (5.446001,-2.200000) (5.452561,-2.160000) (5.459243,-2.120000) (5.466045,-2.080000) (5.472963,-2.040000) (5.479993,-2.000000) (5.487130,-1.960000) (5.494370,-1.920000) (5.501706,-1.880000) (5.509131,-1.840000) (5.516638,-1.800000) (5.524220,-1.760000) (5.531868,-1.720000) (5.539573,-1.680000) (5.547326,-1.640000) (5.555116,-1.600000) (5.562932,-1.560000) (5.570763,-1.520000) (5.578598,-1.480000) (5.586424,-1.440000) (5.594229,-1.400000) (5.602001,-1.360000) (5.609727,-1.320000) (5.617392,-1.280000) (5.624985,-1.240000) (5.632491,-1.200000) (5.639898,-1.160000) (5.647193,-1.120000) (5.654361,-1.080000) (5.661391,-1.040000) (5.668269,-1.000000) (5.674983,-0.960000) (5.681520,-0.920000) (5.687870,-0.880000) (5.694019,-0.840000) (5.699959,-0.800000) (5.705676,-0.760000) (5.711163,-0.720000) (5.716408,-0.680000) (5.721402,-0.640000) (5.726138,-0.600000) (5.730606,-0.560000) (5.734799,-0.520000) (5.738710,-0.480000) (5.742332,-0.440000) (5.745659,-0.400000) (5.748687,-0.360000) (5.751408,-0.320000) (5.753820,-0.280000) (5.755919,-0.240000) (5.757700,-0.200000) (5.759161,-0.160000) (5.760300,-0.120000) (5.761115,-0.080000) (5.761605,-0.040000) (5.761768,0.000000) (5.761605,0.040000) (5.761115,0.080000) (5.760300,0.120000) (5.759161,0.160000) (5.757700,0.200000) (5.755919,0.240000) (5.753820,0.280000) (5.751408,0.320000) (5.748687,0.360000) (5.745659,0.400000) (5.742332,0.440000) (5.738710,0.480000) (5.734799,0.520000) (5.730606,0.560000) (5.726138,0.600000) (5.721402,0.640000) (5.716408,0.680000) (5.711163,0.720000) (5.705676,0.760000) (5.699959,0.800000) (5.694019,0.840000) (5.687870,0.880000) (5.681520,0.920000) (5.674983,0.960000) (5.668269,1.000000) (5.661391,1.040000) (5.654361,1.080000) (5.647193,1.120000) (5.639898,1.160000) (5.632491,1.200000) (5.624985,1.240000) (5.617392,1.280000) (5.609727,1.320000) (5.602001,1.360000) (5.594229,1.400000) (5.586424,1.440000) (5.578598,1.480000) (5.570763,1.520000) (5.562932,1.560000) (5.555116,1.600000) (5.547326,1.640000) (5.539573,1.680000) (5.531868,1.720000) (5.524220,1.760000) (5.516638,1.800000) (5.509131,1.840000) (5.501706,1.880000) (5.494370,1.920000) (5.487130,1.960000) (5.479993,2.000000) (5.472963,2.040000) (5.466045,2.080000) (5.459243,2.120000) (5.452561,2.160000) (5.446001,2.200000) (5.439567,2.240000) (5.433259,2.280000) (5.427080,2.320000) (5.421031,2.360000) (5.415111,2.400000)};
 \node[align=center,anchor=north] at (-3.6,-2.85) {$E=0$ 只要求初始场不变\\这里先加场、后变为理想导体};
 \node[align=center,anchor=north] at (3.6,-2.85) {平衡迈斯纳态，$B\approx0$\\屏蔽电流使磁场绕开样品};
\end{tikzpicture}
\figcap{先加磁场后冷却：理想导体与迈斯纳态的区别（长圆柱截面模型）。}\label{ch26:ref:review-meissner}
\end{center}}

\newcommand{\levitationDiagram}{%
\begin{center}
\begin{tikzpicture}[scale=.88,>=Latex]
 \draw[fill=SoftGray,thick] (-1.5,1.2) rectangle (1.5,2.05);
 \draw[thick] (0,1.2)--(0,2.05);
 \node at (-.75,1.63) {$N$};\node at (.75,1.63) {$S$};
 \draw[fill=PaleBlue,thick] (-2.2,-.6)--(-2.2,-1)
   arc[start angle=180,end angle=360,x radius=2.2,y radius=.5]--(2.2,-.6);
 \draw[fill=PaleBlue,thick] (0,-.6) ellipse (2.2 and .5);
 \node at (0,-.6) {超导圆盘};
 \draw[guide] (1.5,1.63)--(3.8,1.63);
 \fill (4.5,1.63) circle (2.5pt);
 \draw[vector] (4.5,1.63)--(4.5,2.85) node[right] {$F_{\rm m}$};
 \draw[->,TextGray,thick] (4.5,1.63)--(4.5,.41) node[right] {$mg$};
 \node at (4.5,-.1) {磁体受力图};
 \node[align=center,anchor=north] at (0,-2.05) {屏蔽电流改变磁场，磁体受向上磁力\\平衡时 $F_{\rm m}=mg$；稳定性还需考虑几何与钉扎};
\end{tikzpicture}
\figcap{超导磁悬浮的装置位置与磁体受力。}\label{ch26:ref:review-levitation}
\end{center}}

\newcommand{\superconductingRingDiagram}{%
\begin{center}
\begin{tikzpicture}[scale=.82,>=Latex]
 \foreach \cx/\title in {-4/{建立初始状态},0/{缓慢减小外场},4/{外场已撤去}}{
  \draw[DeepBlue,very thick] (\cx,0) circle (1.1);
  \node[align=center] at (\cx,-1.65) {\title};
  \foreach \dx in {-.4,0,.4}{
   \draw[TextGray] (\cx+\dx,0) circle (.09);
   \fill[TextGray] (\cx+\dx,0) circle (.025);
  }
 }
 \foreach \cx in {0,4}{
  \draw[->,DeepBlue,very thick] (\cx+1.4,0)
   arc[start angle=0,end angle=140,radius=1.4];
  \node[DeepBlue,above] at (\cx,1.55) {$I>0$};
 }
 \draw[->,TextGray] (-2.5,0)--(-1.65,0);
 \draw[->,TextGray] (1.7,0)--(2.5,0);
 \node[align=center,anchor=north,font=\small] at (0,-2.2)
 {俯视：$\odot$ 为出纸面正磁通；外磁通减少，逆时针环流产生同向自感磁通\\[5pt]
  左图为初始总场，中图为外场与自感场之和，右图仅为自感场};
\end{tikzpicture}
\figcap{宏观超导环的外场撤去过程。}\label{ch26:ref:review-ring-hold}
\end{center}}

% Historical comparison; values are approximate, not a precision dataset.
\newcommand{\superconductivityChartFaithful}{%
\begin{center}
\resizebox{.98\linewidth}{!}{%
\begin{tikzpicture}[x=.019cm,y=.62cm,>=Latex]
\begin{scope}[shift={(0,0)}]
\draw[axis] (0,.5)--(230,.5) node[right] {$T_c/\mathrm K$};
\foreach \t in {0,50,100,150,200}{
 \draw[TextGray] (\t,.5)--(\t,.65) node[above,font=\scriptsize] {$\t$};
 \draw[TextGray!30,dotted] (\t,.45)--(\t,-20.5);
}
\node[anchor=east,font=\scriptsize,inner sep=2pt] at (0,0) {Hg (1911)};
\draw[DeepBlue,line width=1pt] (0,0)--(4.2,0);
\fill[DeepBlue] (4.2,0) circle (1.8pt);
\node[anchor=west,font=\scriptsize] at (7.2,0) {4.2};
\node[anchor=east,font=\scriptsize,inner sep=2pt] at (0,-1) {Pb (1913)};
\draw[DeepBlue,line width=1pt] (0,-1)--(7.2,-1);
\fill[DeepBlue] (7.2,-1) circle (1.8pt);
\node[anchor=west,font=\scriptsize] at (10.2,-1) {7.2};
\node[anchor=east,font=\scriptsize,inner sep=2pt] at (0,-2) {Nb (1930)};
\draw[DeepBlue,line width=1pt] (0,-2)--(9.2,-2);
\fill[DeepBlue] (9.2,-2) circle (1.8pt);
\node[anchor=west,font=\scriptsize] at (12.2,-2) {9.2};
\node[anchor=east,font=\scriptsize,inner sep=2pt] at (0,-3) {NbN (1941)};
\draw[DeepBlue,line width=1pt] (0,-3)--(16,-3);
\fill[DeepBlue] (16,-3) circle (1.8pt);
\node[anchor=west,font=\scriptsize] at (19,-3) {16};
\node[anchor=east,font=\scriptsize,inner sep=2pt] at (0,-4) {V$_3$Si (1953)};
\draw[DeepBlue,line width=1pt] (0,-4)--(17,-4);
\fill[DeepBlue] (17,-4) circle (1.8pt);
\node[anchor=west,font=\scriptsize] at (20,-4) {17};
\node[anchor=east,font=\scriptsize,inner sep=2pt] at (0,-5) {Nb$_3$Sn (1954)};
\draw[DeepBlue,line width=1pt] (0,-5)--(18,-5);
\fill[DeepBlue] (18,-5) circle (1.8pt);
\node[anchor=west,font=\scriptsize] at (21,-5) {18};
\node[anchor=east,font=\scriptsize,inner sep=2pt] at (0,-6) {Nb$_3$Ge (1973)};
\draw[DeepBlue,line width=1pt] (0,-6)--(23,-6);
\fill[DeepBlue] (23,-6) circle (1.8pt);
\node[anchor=west,font=\scriptsize] at (26,-6) {23};
\node[anchor=east,font=\scriptsize,inner sep=2pt] at (0,-7) {CeCu$_2$Si$_2$ (1979)};
\draw[DeepBlue,line width=1pt] (0,-7)--(0.6,-7);
\fill[DeepBlue] (0.6,-7) circle (1.8pt);
\node[anchor=west,font=\scriptsize] at (3.6,-7) {0.6};
\node[anchor=east,font=\scriptsize,inner sep=2pt] at (0,-8) {UBe$_{13}$ (1983)};
\draw[DeepBlue,line width=1pt] (0,-8)--(0.9,-8);
\fill[DeepBlue] (0.9,-8) circle (1.8pt);
\node[anchor=west,font=\scriptsize] at (3.9,-8) {0.9};
\node[anchor=east,font=\scriptsize,inner sep=2pt] at (0,-9) {UPt$_3$ (1984)};
\draw[DeepBlue,line width=1pt] (0,-9)--(0.5,-9);
\fill[DeepBlue] (0.5,-9) circle (1.8pt);
\node[anchor=west,font=\scriptsize] at (3.5,-9) {0.5};
\node[anchor=east,font=\scriptsize,inner sep=2pt] at (0,-10) {UPd$_2$Al$_3$ (1991)};
\draw[DeepBlue,line width=1pt] (0,-10)--(2,-10);
\fill[DeepBlue] (2,-10) circle (1.8pt);
\node[anchor=west,font=\scriptsize] at (5,-10) {2};
\node[anchor=east,font=\scriptsize,inner sep=2pt] at (0,-11) {CeCoIn$_5$ (2001)};
\draw[DeepBlue,line width=1pt] (0,-11)--(2.3,-11);
\fill[DeepBlue] (2.3,-11) circle (1.8pt);
\node[anchor=west,font=\scriptsize] at (5.3,-11) {2.3};
\node[anchor=east,font=\scriptsize,inner sep=2pt] at (0,-12) {PuCoGa$_5$ (2002)};
\draw[DeepBlue,line width=1pt] (0,-12)--(18.5,-12);
\fill[DeepBlue] (18.5,-12) circle (1.8pt);
\node[anchor=west,font=\scriptsize] at (21.5,-12) {18.5};
\node[anchor=east,font=\scriptsize,inner sep=2pt] at (0,-13) {PuRhGa$_5$ (2004)};
\draw[DeepBlue,line width=1pt] (0,-13)--(9,-13);
\fill[DeepBlue] (9,-13) circle (1.8pt);
\node[anchor=west,font=\scriptsize] at (12,-13) {9};
\node[anchor=east,font=\scriptsize,inner sep=2pt] at (0,-14) {LaBaCuO (1986)};
\draw[DeepBlue,line width=1pt] (0,-14)--(30,-14);
\fill[DeepBlue] (30,-14) circle (1.8pt);
\node[anchor=west,font=\scriptsize] at (33,-14) {30};
\node[anchor=east,font=\scriptsize,inner sep=2pt] at (0,-15) {LaSrCuO (1987)};
\draw[DeepBlue,line width=1pt] (0,-15)--(36,-15);
\fill[DeepBlue] (36,-15) circle (1.8pt);
\node[anchor=west,font=\scriptsize] at (39,-15) {36};
\node[anchor=east,font=\scriptsize,inner sep=2pt] at (0,-16) {YBaCuO (1987)};
\draw[DeepBlue,line width=1pt] (0,-16)--(93,-16);
\fill[DeepBlue] (93,-16) circle (1.8pt);
\node[anchor=west,font=\scriptsize] at (96,-16) {93};
\node[anchor=east,font=\scriptsize,inner sep=2pt] at (0,-17) {BiSrCaCuO (1988)};
\draw[DeepBlue,line width=1pt] (0,-17)--(110,-17);
\fill[DeepBlue] (110,-17) circle (1.8pt);
\node[anchor=west,font=\scriptsize] at (113,-17) {110};
\node[anchor=east,font=\scriptsize,inner sep=2pt] at (0,-18) {TlBaCaCuO (1988)};
\draw[DeepBlue,line width=1pt] (0,-18)--(125,-18);
\fill[DeepBlue] (125,-18) circle (1.8pt);
\node[anchor=west,font=\scriptsize] at (128,-18) {125};
\node[anchor=east,font=\scriptsize,inner sep=2pt] at (0,-19) {HgBaCaCuO (1993)};
\draw[DeepBlue,line width=1pt] (0,-19)--(133,-19);
\fill[DeepBlue] (133,-19) circle (1.8pt);
\node[anchor=west,font=\scriptsize] at (136,-19) {133};
\node[anchor=east,font=\scriptsize,inner sep=2pt] at (0,-20) {HgBaCaCuO，30 GPa};
\draw[DeepBlue,line width=1pt] (0,-20)--(164,-20);
\fill[DeepBlue] (164,-20) circle (1.8pt);
\node[anchor=west,font=\scriptsize] at (167,-20) {164};
\end{scope}
\begin{scope}[shift={(470,0)}]
\draw[axis] (0,.5)--(230,.5) node[right] {$T_c/\mathrm K$};
\foreach \t in {0,50,100,150,200}{
 \draw[TextGray] (\t,.5)--(\t,.65) node[above,font=\scriptsize] {$\t$};
 \draw[TextGray!30,dotted] (\t,.45)--(\t,-20.5);
}
\node[anchor=east,font=\scriptsize,inner sep=2pt] at (0,0) {HgTlBaCaCuO (1995)};
\draw[DeepBlue,line width=1pt] (0,0)--(138,0);
\fill[DeepBlue] (138,0) circle (1.8pt);
\node[anchor=west,font=\scriptsize] at (141,0) {138};
\node[anchor=east,font=\scriptsize,inner sep=2pt] at (0,-1) {BKBO (1988)};
\draw[DeepBlue,line width=1pt] (0,-1)--(30,-1);
\fill[DeepBlue] (30,-1) circle (1.8pt);
\node[anchor=west,font=\scriptsize] at (33,-1) {30};
\node[anchor=east,font=\scriptsize,inner sep=2pt] at (0,-2) {YPd$_2$B$_2$C (1994)};
\draw[DeepBlue,line width=1pt] (0,-2)--(23,-2);
\fill[DeepBlue] (23,-2) circle (1.8pt);
\node[anchor=west,font=\scriptsize] at (26,-2) {23};
\node[anchor=east,font=\scriptsize,inner sep=2pt] at (0,-3) {MgB$_2$ (2001)};
\draw[DeepBlue,line width=1pt] (0,-3)--(39,-3);
\fill[DeepBlue] (39,-3) circle (1.8pt);
\node[anchor=west,font=\scriptsize] at (42,-3) {39};
\node[anchor=east,font=\scriptsize,inner sep=2pt] at (0,-4) {Li，33 GPa (2002)};
\draw[DeepBlue,line width=1pt] (0,-4)--(16,-4);
\fill[DeepBlue] (16,-4) circle (1.8pt);
\node[anchor=west,font=\scriptsize] at (19,-4) {16};
\node[anchor=east,font=\scriptsize,inner sep=2pt] at (0,-5) {K$_3$C$_{60}$ (1991)};
\draw[DeepBlue,line width=1pt] (0,-5)--(18,-5);
\fill[DeepBlue] (18,-5) circle (1.8pt);
\node[anchor=west,font=\scriptsize] at (21,-5) {18};
\node[anchor=east,font=\scriptsize,inner sep=2pt] at (0,-6) {RbCsC$_{60}$ (1991)};
\draw[DeepBlue,line width=1pt] (0,-6)--(33,-6);
\fill[DeepBlue] (33,-6) circle (1.8pt);
\node[anchor=west,font=\scriptsize] at (36,-6) {33};
\node[anchor=east,font=\scriptsize,inner sep=2pt] at (0,-7) {Cs$_3$C$_{60}$，1.4 GPa};
\draw[DeepBlue,line width=1pt] (0,-7)--(38,-7);
\fill[DeepBlue] (38,-7) circle (1.8pt);
\node[anchor=west,font=\scriptsize] at (41,-7) {38};
\node[anchor=east,font=\scriptsize,inner sep=2pt] at (0,-8) {CNT低温报道 (2001)};
\draw[DeepBlue,line width=1pt] (0,-8)--(0.55,-8);
\fill[DeepBlue] (0.55,-8) circle (1.8pt);
\node[anchor=west,font=\scriptsize] at (3.55,-8) {0.55};
\node[anchor=east,font=\scriptsize,inner sep=2pt] at (0,-9) {CNT高温报道 (2006)};
\draw[DeepBlue,line width=1pt] (0,-9)--(12,-9);
\fill[DeepBlue] (12,-9) circle (1.8pt);
\node[anchor=west,font=\scriptsize] at (15,-9) {12};
\node[anchor=east,font=\scriptsize,inner sep=2pt] at (0,-10) {掺硼金刚石 (2004)};
\draw[DeepBlue,line width=1pt] (0,-10)--(4,-10);
\fill[DeepBlue] (4,-10) circle (1.8pt);
\node[anchor=west,font=\scriptsize] at (7,-10) {4};
\node[anchor=east,font=\scriptsize,inner sep=2pt] at (0,-11) {LaOFeP (2006)};
\draw[DeepBlue,line width=1pt] (0,-11)--(4,-11);
\fill[DeepBlue] (4,-11) circle (1.8pt);
\node[anchor=west,font=\scriptsize] at (7,-11) {4};
\node[anchor=east,font=\scriptsize,inner sep=2pt] at (0,-12) {LaOFFeAs (2008)};
\draw[DeepBlue,line width=1pt] (0,-12)--(26,-12);
\fill[DeepBlue] (26,-12) circle (1.8pt);
\node[anchor=west,font=\scriptsize] at (29,-12) {26};
\node[anchor=east,font=\scriptsize,inner sep=2pt] at (0,-13) {铁基材料高温区 (2008)};
\draw[DeepBlue,line width=1pt] (0,-13)--(43,-13);
\fill[DeepBlue] (43,-13) circle (1.8pt);
\node[anchor=west,font=\scriptsize] at (46,-13) {43};
\node[anchor=east,font=\scriptsize,inner sep=2pt] at (0,-14) {掺杂SrFFeAs (2009)};
\draw[DeepBlue,line width=1pt] (0,-14)--(56,-14);
\fill[DeepBlue] (56,-14) circle (1.8pt);
\node[anchor=west,font=\scriptsize] at (59,-14) {56};
\node[anchor=east,font=\scriptsize,inner sep=2pt] at (0,-15) {FeSe单层报道 (2015)};
\draw[DeepBlue,line width=1pt] (0,-15)--(109,-15);
\fill[DeepBlue] (109,-15) circle (1.8pt);
\node[anchor=west,font=\scriptsize] at (112,-15) {109};
\node[anchor=east,font=\scriptsize,inner sep=2pt] at (0,-16) {硫氢化物，155 GPa};
\draw[DeepBlue,line width=1pt] (0,-16)--(203,-16);
\fill[DeepBlue] (203,-16) circle (1.8pt);
\node[anchor=west,font=\scriptsize] at (206,-16) {203};
\node[anchor=east,font=\scriptsize,inner sep=2pt] at (0,-17) {NbC（补充）};
\draw[DeepBlue,line width=1pt] (0,-17)--(11,-17);
\fill[DeepBlue] (11,-17) circle (1.8pt);
\node[anchor=west,font=\scriptsize] at (14,-17) {11};
\node[anchor=east,font=\scriptsize,inner sep=2pt] at (0,-18) {Nb$_3$Al（补充）};
\draw[DeepBlue,line width=1pt] (0,-18)--(18,-18);
\fill[DeepBlue] (18,-18) circle (1.8pt);
\node[anchor=west,font=\scriptsize] at (21,-18) {18};
\node[anchor=east,font=\scriptsize,inner sep=2pt] at (0,-19) {CaC$_6$（补充）};
\draw[DeepBlue,line width=1pt] (0,-19)--(11.5,-19);
\fill[DeepBlue] (11.5,-19) circle (1.8pt);
\node[anchor=west,font=\scriptsize] at (14.5,-19) {11.5};
\end{scope}
\end{tikzpicture}}
\figcap{不同超导材料族的历史温区。年代与温度为约值；高压、掺杂和单层条件不能忽略，CNT与单层FeSe列的是历史报道温区，不能直接等同为常压块材零电阻温度。}\label{ch26:fig:s41-10-9}
\end{center}
液态冷媒提供的温区约为：液氦 $4.2\,\mathrm K$、液氢 $20\,\mathrm K$、液氮 $77\,\mathrm K$、液态四氟化碳 $145\,\mathrm K$。
铜氧化物和铁基材料的简式表示材料族，不表示未经掺杂的任意同名样品都具有图示 $T_c$。
图中 MgB$_2$、富勒烯、重费米子、碳材料及铁基体系分布于不同温区，说明新材料的发现并不总是刷新最高温度纪录；高压硫氢化物还需与常压材料分开评价。
}


电与磁最初分别由静电学、稳恒电流和电磁感应来描述。图~\ref{ch26:ref:review-maxwell-portrait} 中的麦克斯韦（James Clerk Maxwell）的关键工作，是看出这些规律必须组成
一个彼此相容的整体；而卡末林·昂内斯（Heike Kamerlingh Onnes）在低温实验中发现的超导现象，则把电磁场与物质的相互作用推到了一个
极端而清晰的情形：电阻可以消失，磁通却不能只按“零电阻导体”来理解。图~\ref{ch26:ref:review-history-context} 将这两条历史线索并列，左侧是场的理论，右侧是场与低温物质的实验。

电磁感应定律揭示了变化的磁场能够产生电场。由此自然会追问：变化的电场能否产生磁场？若原有的
安培定律只能处理稳恒电流，它在含有随时间变化的电荷与电场时是否仍然完备？麦克斯韦在尚无直接
实验依据的情况下，从理论自洽性出发补上位移电流项，使电场和磁场能够彼此激发。由此得到的电磁波
传播速度与当时测得的光速一致，他进一步提出光的电磁理论，经典电磁学的基本框架遂告建立。

这一过程值得特别注意：麦克斯韦并不是把彼此无关的四条实验定律简单抄在一起，而是要求它们同时
满足电荷守恒、空间对称性和有限传播速度。静电场的高斯定律告诉我们电荷能够成为电场的源，磁场
高斯定律则表明磁感线没有已经被发现的起点和终点；法拉第定律说明变化磁场会在周围激起涡旋电场。
唯独旧的安培定律只容纳稳恒传导电流，在时间变化问题中显得不对称。位移电流修正消除了这种
不对称，并使方程组能够在没有电荷、没有导线的真空中仍然具有随时间传播的非零解。电磁波并非额外
塞进理论的假设，而是完备方程组自身推出来的结果。

从学习方法上说，本章的难点不在记住更多公式，而在分清三种层次。第一层是实验事实，例如电荷产生
电场、电流产生磁场以及电磁感应；第二层是把事实压缩为局部方程，并检查这些方程能否同时成立；第三层
才是从方程推出波速、能流和辐射。只记“变化电场产生磁场”会漏掉为什么必须有这一项；只记四条方程
又会看不出电磁波从何而来。后文的每一步都围绕“源—场—传播—能量”这条线索展开。

\maxwellOnnesContext

\historyfigure[0.27\linewidth]{maxwell-built-in-v1.png}%
{麦克斯韦（James Clerk Maxwell，1831--1879），历史人物肖像。\label{ch26:ref:review-maxwell-portrait}}\label{ch26:fig:s41-10-1}


\subsection{位移电流：补上安培定律的缺口}


本段限定真空或磁导率近似为 $\mu_0$ 的非磁介质，$\vect j$ 为自由传导电流密度，$\rho$ 为自由电荷密度。含磁化介质时应使用 $\vect H$，不能仍把任意介质中的 $\vect B$ 直接除以 $\mu_0$。
稳恒电流产生磁场时，安培环路定律为
\begin{equation}
  \oint_C\vect B\cdot\dd\vect l=\mu_0 I_{\rm c},
\end{equation}
其中 $I_{\rm c}=\int_S\vect j\cdot\dd\vect S$ 是穿过跨越面 $S$ 的传导电流。若电流处处稳恒，选择
同一边界 $C$ 的不同跨越面，所得 $I_{\rm c}$ 相同，公式没有歧义。

先看第一个矛盾。设一块均匀、弱导电的球形介质中心原有电荷，电荷随时间沿径向向外泄漏。
电流密度具有严格的球对称性。若在球内任取一条闭合环路，按原来的安培环路定律，只要跨越面上
有净电流穿过，环路周围似乎就应当出现环向磁场。然而球对称系统中不存在一个在每一点都非零、又能
选出固定切向的矢量场；任何非零环向磁场都会破坏球对称性。因此，单独用传导电流作磁场源会导出
相互冲突的结论。矛盾的根源是：电荷流出时，球内电荷量和由高斯定律决定的电场也同时在变化。

把这个论证写得更定量一些。若半径 $r$ 的球面内电荷为 $Q_r(t)$，球对称性给出
\begin{equation}
  4\pi r^2D_r=Q_r(t).
\end{equation}
设向外传导电流为 $I_r$，电荷守恒要求 $I_r=-\dd Q_r/\dd t$；另一方面，穿过同一球面的位移电流为
\begin{equation}
  I_{\rm d}=\frac{\dd}{\dd t}(4\pi r^2D_r)=\frac{\dd Q_r}{\dd t}=-I_r.
\end{equation}
球对称性还把这个积分结论加强为每一点的关系：
\[
 j_r=\frac{I_r}{4\pi r^2},\qquad
 \frac{\partial D_r}{\partial t}=-\frac{I_r}{4\pi r^2},
 \qquad \vect j+\frac{\partial\vect D}{\partial t}=0.
\]
因此不是只在闭合面上“总量抵消”，而是局部的两种源也相消。
无外加磁场时，对称性只允许径向磁场；再由任意同心球面磁通为零，可得 $B_r=0$。于是球对称解确实为 $\vect B=0$。

图~\ref{ch26:fig:s41-10-2} 中实线箭头表示径向电场，灰色虚线箭头表示径向传导电流。任选闭合环路 $\Gamma$ 都不能在
球对称条件下挑出唯一的环向磁场；图右把同一点的 $\vect j$ 与 $\partial\vect D/\partial t$ 单独放大，两者大小相等、方向相反。

\radialCurrentSphereDiagram

给电容器充电时，问题暴露出来：取 $S_1$ 穿过导线，有 $I_{\rm c}=I$；取 $S_2$ 从两极板之间鼓出，
它不切断任何导线，因而 $I_{\rm c}=0$。两个面有同一边界，却会给出两个不同的磁场环量。

这里真正参与定律的是“同一闭合回路的任意跨越面”。环路 $C$ 可以想成一根铁丝围成的圈，$S_1$ 与
$S_2$ 则像两张边缘都粘在铁丝上的橡皮膜。一张膜穿过导线，另一张膜向电容器间隙鼓起；二者形状不同，
边界却完全相同。若右端只计算传导电流，左端同一个 $\oint_C\vect B\cdot\dd\vect l$ 就会得到两个值，
这不是近似误差，而是定律内部失去定义。位移电流项的任务正是让右端也只依赖边界，而不依赖所选曲面。

图~\ref{ch26:ref:review-surface-choice} 中两种蓝色剖面具有相同边界；从左图变为右图时，跨越面绕开左极板而不切断导线。

\displacementCurrentDiagram

极板间虽然没有自由电荷穿越，电位移通量却在改变。麦克斯韦定义位移电流
\begin{equation}
  I_{\rm d}=\frac{\dd}{\dd t}\int_S\vect D\cdot\dd\vect S,
\end{equation}
一般介质中的形式为 $\oint_C\vect H\cdot\dd\vect l=\int_S(\vect j+\partial\vect D/\partial t)\cdot\dd\vect S$。在本段的非磁介质条件下，它写为
\begin{equation}
  \boxed{\oint_C\vect B\cdot\dd\vect l
  =\mu_0\left(\int_S\vect j\cdot\dd\vect S
  +\frac{\dd}{\dd t}\int_S\vect D\cdot\dd\vect S\right)}.
\end{equation}
真空中 $\vect D=\varepsilon_0\vect E$。平行板电容器忽略边缘效应时，板间
$E=Q/(\varepsilon_0 A)$，于是
\begin{equation}
  I_{\rm d}=\varepsilon_0\frac{\dd(EA)}{\dd t}
  =\frac{\dd Q}{\dd t}=I.
\end{equation}
因此 $S_1$ 上的传导电流与 $S_2$ 上的位移电流恰好相等，磁场环量不再依赖跨越面的选取。

若在极板之间取半径为 $r$、与极板同轴的圆形环路，并假设 $r$ 小于极板半径，则该环路所围面积中的
位移电流为 $I_{{\rm d},r}=I\pi r^2/A$。由圆柱对称性，环路上的磁感应强度大小处处相同且沿切向，故
\begin{equation}
  B(2\pi r)=\mu_0I\frac{\pi r^2}{A},
  \qquad B=\frac{\mu_0Ir}{2A}.
\end{equation}
在极板投影之外，所包围的总位移电流已等于 $I$，于是 $B=\mu_0I/(2\pi r)$。这表明充电间隙中的
磁场并非抽象的“修补量”：它同样可由环路定律给出具体空间分布，并与导线周围磁场连续衔接。

回看弱导电球的例子，穿过闭合面的向外传导电流使内部电荷减少，与此同时电通量也减少。传导电流与
位移电流的代数和满足连续性，因而不会再凭空选出一个破坏球对称性的磁场方向。两个表面上看不同的
矛盾，其共同点正是都遗漏了变化电场的磁效应。

这个修正并非只为解决电容器的特殊矛盾。局部电荷守恒要求
\begin{equation}
  \nabla\cdot\vect j=-\frac{\partial\rho}{\partial t}.
\end{equation}
若仍用 $\nabla\times\vect B=\mu_0\vect j$，两边取散度便得到
$0=\mu_0\nabla\cdot\vect j$，它只允许电荷密度不随时间变化。加入位移电流后，利用
$\nabla\cdot\vect D=\rho$，有
\begin{equation}
  0=\mu_0\left(\nabla\cdot\vect j+
  \frac{\partial}{\partial t}\nabla\cdot\vect D\right)
  =\mu_0\left(\nabla\cdot\vect j+\frac{\partial\rho}{\partial t}\right),
\end{equation}
恰好还原连续性方程。

图~\ref{ch26:ref:review-continuity} 从上向下补入位移电流后，散度方程恢复一般的电荷守恒。

\continuityDiagram

这里的“位移电流”不是极板间真的有电荷横渡真空；它是变化电场对磁场的贡献，单位与电流相同。
这个名称带有明显的历史色彩。麦克斯韦所处的时代仍常把电磁场设想为某种介质的状态，认为外场会使
假想的真空介质或真实电介质发生极化，极化状态的变化仿佛形成一种电流。今天已经不再需要这种“以太”
图景，但“位移电流”这一名称以及电位移矢量 $\vect D$ 被保留下来。现代观点强调的不是机械介质真的
流动，而是 $\partial\vect D/\partial t$ 与传导电流一样，是涡旋磁场的源。

位移电流与传导电流还应在物理机制上区分。金属导线中的传导电流伴随自由载流子的定向运动，可以在
截面两侧统计净通过的电荷；真空电容间隙中的位移电流则没有这样的电荷输运。二者之所以能写在同一个
括号中，是因为它们对磁场环量产生同样的源效应，并共同保证总电流连续。不能据此断言“真空中也有
电子从一块极板跳到另一块极板”，也不能把位移电流当作某种额外漏电。


% EM teaching enrichment 20260927
\paragraph{位移电流不是电子穿过真空间隙}
平行板充电时，导线中传导电流改变极板自由电荷，板间的变化电场对应位移电流。二者在磁场环量关系中以相同单位出现，却描述不同机制。
对于均匀线性介质，适用的宏观项为 $\partial\boldsymbol D/\partial t$，包括真空项与极化变化；若介质还漏电，则传导项也要保留。
检查一个跨越面时应分别问：有没有自由电荷穿面？穿面的电位移通量是否变化？只凭“这里没有导线”不能把安培--麦克斯韦方程的右边判为零。


\subsection{麦克斯韦方程组}


在真空的宏观描述中，电磁场的四条基本方程可写成
\begin{align}
  \nabla\cdot\vect E&=\frac{\rho}{\varepsilon_0},
  &\oint_{\partial V}\vect E\cdot\dd\vect A&=\frac{Q_V}{\varepsilon_0},
  \\
  \nabla\cdot\vect B&=0,
  &\oint_{\partial V}\vect B\cdot\dd\vect A&=0,
  \\
  \nabla\times\vect E&=-\frac{\partial\vect B}{\partial t},
  &\oint_C\vect E\cdot\dd\vect l&=-\frac{\dd}{\dd t}\int_S\vect B\cdot\dd\vect S,
  \\
  \nabla\times\vect B&=\mu_0\vect j+\mu_0\varepsilon_0
  \frac{\partial\vect E}{\partial t},
  &\oint_C\vect B\cdot\dd\vect l&=\mu_0 I_{\rm c}+\mu_0\varepsilon_0
  \frac{\dd}{\dd t}\int_S\vect E\cdot\dd\vect S.
\end{align}
左列是微分形式，描述一个点附近的场；右列是积分形式，直接联系闭合曲面或闭合回路上的整体量。
四条方程分别表达：电荷是电场的源；没有观测到孤立磁荷；变化的磁场产生涡旋电场；传导电流与
变化的电场共同产生涡旋磁场。

积分形式适合直接分析具有对称性的有限区域，微分形式则把定律落实到空间中的每一点。例如
$\nabla\cdot\vect E=\rho/\varepsilon_0$ 在直角坐标系中就是
\begin{equation}
  \frac{\partial E_x}{\partial x}+\frac{\partial E_y}{\partial y}
  +\frac{\partial E_z}{\partial z}=\frac{\rho}{\varepsilon_0},
\end{equation}
而法拉第定律包含三个按 $x\to y\to z\to x$ 循环轮换得到的分量式，例如
\begin{equation}
  \frac{\partial E_z}{\partial y}-\frac{\partial E_y}{\partial z}
  =-\frac{\partial B_x}{\partial t}.
\end{equation}
安培--麦克斯韦定律同样给出三个分量关系，例如
\begin{equation}
  \frac{\partial B_z}{\partial y}-\frac{\partial B_y}{\partial z}
  =\mu_0j_x+\mu_0\varepsilon_0\frac{\partial E_x}{\partial t}.
\end{equation}
这些式子可由积分定律分别作用在无限小的高斯盒或矩形回路上得到。中学竞赛很少要求逐分量求解
完整方程组，但看清分量间的循环结构，有助于理解后面的平面波推导不是另立的新规律。

完整写出两个旋度方程的分量结构，可以看到这种循环关系：
\begin{align}
 \frac{\partial E_z}{\partial y}-\frac{\partial E_y}{\partial z}
   &=-\frac{\partial B_x}{\partial t},\\
 \frac{\partial E_x}{\partial z}-\frac{\partial E_z}{\partial x}
   &=-\frac{\partial B_y}{\partial t},\\
 \frac{\partial E_y}{\partial x}-\frac{\partial E_x}{\partial y}
   &=-\frac{\partial B_z}{\partial t},\\
 \frac{\partial B_z}{\partial y}-\frac{\partial B_y}{\partial z}
   &=\mu_0j_x+\mu_0\varepsilon_0\frac{\partial E_x}{\partial t},\\
 \frac{\partial B_x}{\partial z}-\frac{\partial B_z}{\partial x}
   &=\mu_0j_y+\mu_0\varepsilon_0\frac{\partial E_y}{\partial t},\\
 \frac{\partial B_y}{\partial x}-\frac{\partial B_x}{\partial y}
   &=\mu_0j_z+\mu_0\varepsilon_0\frac{\partial E_z}{\partial t}.
\end{align}
每一个旋度分量都描述“垂直于某一小面的环流密度”。这也解释了为什么平面波中只需选择一对互相
垂直的电、磁分量，另外几条分量式就自动化为 $0=0$。

方程组还需要洛伦兹力定律
\begin{equation}
  \boxed{\vect F=q(\vect E+\vect v\times\vect B)}
\end{equation}
才能说明场如何作用于带电粒子。麦克斯韦方程决定“场怎样产生和传播”，洛伦兹力决定“场怎样推动
物质”，二者分工不同。

还应避免把“场的方程”和“粒子的运动方程”混为一谈。给定 $\rho,\vect j$ 与边界条件，麦克斯韦方程
求出 $\vect E,\vect B$；再把所得场代入洛伦兹力，结合牛顿定律，才能求粒子的轨迹。反过来，
粒子运动又会改变电荷、电流分布并重新产生场。完整的电磁动力学是这种相互耦合，而不是先验地把场
当成不会受带电粒子影响的固定背景。

图~\ref{ch26:ref:review-maxwell-light} 中连接三个方框的箭头分别表示联立求解和比较传播速度，不是对实验定律的简单罗列。

\maxwellLightPanel


% EM teaching enrichment 20260927
\paragraph{四条方程与材料关系各管什么}
电的通量方程约束电荷源，磁的通量方程排除磁单极源；两条旋度方程描述变化场及电流产生的环流。它们需与电荷运动方程、初值、边界条件共同使用，不是四个互不相关的算式。
例如对安培--麦克斯韦方程取散度，左侧为零，代入电的高斯定律后得到
$\partial\rho/\partial t+\nabla\cdot\boldsymbol J=0$。若删去位移电流项，就会错误地强迫任何问题都满足 $\nabla\cdot\boldsymbol J=0$，无法描述极板正在积累电荷的过程。
介质中还需给出 $\boldsymbol D(\boldsymbol E)$、$\boldsymbol B(\boldsymbol H)$ 等材料关系；这些关系不是所有材料都相同的基本场方程。


\subsection{位移电流的定量应用}


真空中的经典电磁场满足
\begin{align}
 \nabla\cdot\vect E&=\frac{\rho}{\varepsilon_0},&
 \nabla\cdot\vect B&=0,\\
 \nabla\times\vect E&=-\frac{\partial\vect B}{\partial t},&
 \nabla\times\vect B&=\mu_0\vect J+\mu_0\varepsilon_0
 \frac{\partial\vect E}{\partial t}.
\end{align}
最后一式中的 $\varepsilon_0\partial\vect E/\partial t$ 是位移电流密度。它保证电荷守恒与安培环路定律相容：
对充电电容器，无论选取跨过导线还是穿过极板间隙的曲面，环路积分都给出同一磁场。


下面从充电电流出发，确定极板间变化的电场如何产生磁场。

\begin{origexamplebox}{平行板电容器中的位移电流}
真空圆形平行板电容器极板半径为 $R$，板间距远小于 $R$，在拟稳近似下忽略边缘效应，以恒定电流 $I$ 充电。沿轴线观察时，板间电场沿出纸面方向增大，如图~\ref{ch26:fig:displacement-question}。求板间距轴线 $r<R$ 处的磁感应强度，并说明径向跨过 $r=R$ 时结果如何衔接。
% Question contract: show plate radius and observation point only; B and the 安培 loop belong in the solution.
\begin{center}
\begin{tikzpicture}[scale=.9,>=Latex]
 \draw[very thick] (0,0) circle (1.8);
 \foreach \x/\y in {0/0,.55/.55,-.55/.55,.55/-.55,-.55/-.55,1.35/0,-1.35/0,0/1.35,0/-1.35}{
   \draw[TextGray] (\x,\y) circle (.07);\fill[TextGray] (\x,\y) circle (.018);
 }
 \draw[TextGray] (0,0)--(-1.485,1.018) node[above left] {$R$};
 \draw[TextGray] (0,0)--(.743,-.743) node[midway,below left] {$r$};
 \fill[DeepBlue] (.743,-.743) circle (2pt) node[right] {$P$};
 \node[below] at (0,-2) {沿板间轴线观察};
 \node[align=left,anchor=west] at (2.3,.3)
  {$\odot$：电场沿出纸面方向增大\\[4pt]
   圆周：极板边缘的投影\\[4pt]
   $P$ 点的磁场待求};
\end{tikzpicture}
\figcap{极板半径、观察点位置及电场变化方向。}\label{ch26:fig:displacement-question}
\end{center}
\end{origexamplebox}
\begin{xesolution}
图~\ref{ch26:fig:displacement-current} 给出解答中选取的安培环路及相应磁场方向。
\begin{center}
\begin{tikzpicture}[scale=.9,>=Latex]
 \draw[very thick] (0,0) circle (1.8);
 \draw[DeepBlue,dashed,thick] (0,0) circle (1.05);
 \foreach \x/\y in {0/0,.55/.55,-.55/.55,.55/-.55,-.55/-.55,1.35/0,-1.35/0,0/1.35,0/-1.35}{
   \draw[TextGray] (\x,\y) circle (.07);\fill[TextGray] (\x,\y) circle (.018);
 }
 \draw[->,DeepBlue,thick] (1.05,0) arc[start angle=0,end angle=70,radius=1.05];
 \draw[TextGray] (0,0)--(-1.485,1.018);
 \node[above left] at (-1.485,1.018) {$R$};
 \draw[TextGray] (0,0)--(.743,-.743);
 \node[below left] at (.44,-.38) {$r$};
 \node[right,DeepBlue] at (1.08,.72) {$\vect B$};
 \node[below] at (0,-2) {沿板间轴线观察};
 \node[align=left,anchor=west] at (2.7,.5)
  {$\odot$：$\partial\vect E/\partial t$ 出纸面\\[4pt]
   虚线圆：安培环路 $C$\\[4pt]
   实线圆：极板边缘的投影\\[4pt]
   $I_{\rm d}(r)=I\,r^2/R^2$};
\end{tikzpicture}
\figcap{充电电容器中的位移电流。环路 $C$ 围成的曲面穿过极板间隙时，由变化电场补足传导电流。}\label{ch26:fig:displacement-current}
\end{center}

由 $Q=\varepsilon_0E\pi R^2$ 和 $\dd Q/\dd t=I$，有 $\varepsilon_0\partial E/\partial t=I/(\pi R^2)$。板间电场近似均匀，总位移电流等于传导电流 $I$。半径为 $r$ 的安培环路所包围的电场面积占总极板面积的 $r^2/R^2$，故
\[
 I_d(r)=I\frac{r^2}{R^2}.
\]
由安培--麦克斯韦定律
\[
 B(2\pi r)=\mu_0I_d(r),
 \qquad
 \boxed{B(r)=\frac{\mu_0Ir}{2\pi R^2}\quad(r<R)}.
\]
在同一横截面上，$r\ge R$ 的环路包围全部位移电流，故 $B(r)=\mu_0I/(2\pi r)$。两支表达式在 $r=R$ 相等，验证径向连续性；这不是把板间观察点直接移到任意一段真实导线附近。
图中的位移电流出纸面，右手拇指朝向读者，四指给出逆时针磁场方向。$r\to0$ 时 $B\to0$，轴线上不应存在优选切向；$\mu_0Ir/R^2$ 的单位为特斯拉。
\end{xesolution}


% EM structure and teaching revision 20260928
\subsection{用连续性检验位移电流项}
对安培--麦克斯韦方程取散度，利用旋度的散度为零，可得
\[
 0=\nabla\cdot\boldsymbol j+\frac{\partial}{\partial t}(\nabla\cdot\boldsymbol D)
 =\nabla\cdot\boldsymbol j+\frac{\partial\rho_f}{\partial t}.
\]
这正是自由电荷的连续性方程。位移电流不是为了给真空人为安排真实载流子，而是使变化电场、磁场与电荷守恒能够在同一组场方程中成立。
对正在充电的理想平行板电容器，取同一条闭合环路，可以选穿过导线的曲面，也可以选绕开导线、穿过极板间隙的曲面。前者包围传导电流，后者包围电位移通量的变化率，两者应给出相同环量。
若只取半径 $r$ 小于极板半径 $a$ 的同心圆回路，并忽略边缘场，极板间 $D$ 均匀，故包围的位移电流为总充电电流的 $r^2/a^2$。磁场因而随 $r$ 线性增加，而不是处处按整根导线的 $1/r$ 规律变化。曲面选择与包围比例，是这一类题目最容易混淆的两个步骤。


\section{电磁波、能流与辐射压}

\subsection{真空中的电磁波}


在没有电荷和传导电流的真空区域，$\rho=0,\vect j=0$。设一列平面波沿 $+x$ 方向传播，场具有共同相位
\begin{equation}
  \vect E=\vect E_0\cos(kx-\omega t),
  \qquad
  \vect B=\vect B_0\cos(kx-\omega t).
\end{equation}
由 $\nabla\cdot\vect E=0$ 得 $-kE_{0x}\sin(kx-\omega t)=0$，故 $E_{0x}=0$；同理由
$\nabla\cdot\vect B=0$ 得 $B_{0x}=0$。因此源自由空间中的平面电磁波是横波。

不妨选 $\vect E=E_0\cos(kx-\omega t)\uvec y$，并设
$\vect B=B_0\cos(kx-\omega t)\uvec z$。法拉第定律的 $z$ 分量给出
\begin{equation}
  \frac{\partial E_y}{\partial x}=-\frac{\partial B_z}{\partial t}
  \quad\Longrightarrow\quad
  kE_0=\omega B_0.
\end{equation}
安培--麦克斯韦定律的 $y$ 分量为
\begin{equation}
  -\frac{\partial B_z}{\partial x}
  =\mu_0\varepsilon_0\frac{\partial E_y}{\partial t}
  \quad\Longrightarrow\quad
  kB_0=\mu_0\varepsilon_0\omega E_0.
\end{equation}
消去振幅便得到
\begin{equation}
  \boxed{\frac{\omega}{k}=\frac{1}{\sqrt{\mu_0\varepsilon_0}}\equiv c},
  \qquad
  \boxed{E_0=cB_0}.
\end{equation}
这个由电磁常量预言的速度等于真空光速，是“光就是电磁波”的核心证据。更一般地，三者满足
\begin{equation}
  \boxed{\vect B=\frac{1}{\omega}\vect k\times\vect E},
  \qquad
  \vect E\perp\vect B,\quad \vect E\perp\vect k,
  \quad \vect B\perp\vect k.
\end{equation}

推导中场的振幅 $E_0,B_0$ 在方程两边被约去，这并不意味着振幅为零，而是说明真空常量只决定传播
速度及两种场的振幅比，实际振幅由波源强弱决定。频率同样由波源的振荡频率决定；一旦频率 $f$ 给定，
真空中的波长便不能任意选择，而必须满足
\begin{equation}
  \lambda=\frac{c}{f},\qquad k=\frac{2\pi}{\lambda},\qquad \omega=2\pi f.
\end{equation}
传播方向的垂直平面有两个彼此独立的偏振方向。任意线偏振可以分解成这两个方向的叠加；若两分量有
相位差，则可形成椭圆偏振或圆偏振。对这里的真空平面波，$\vect E$ 与 $\vect B$ 同相，三者按
$\vect E,\vect B,\vect k$ 的次序构成右手系。

也可以不先猜传播速度，而直接从麦克斯韦方程消去磁场。对法拉第定律取旋度，使用恒等式
$\nabla\times(\nabla\times\vect E)=\nabla(\nabla\cdot\vect E)-\nabla^2\vect E$；在真空无源区
$\nabla\cdot\vect E=0$，再代入安培--麦克斯韦定律，得到
\begin{equation}
  \nabla^2\vect E-\mu_0\varepsilon_0\frac{\partial^2\vect E}{\partial t^2}=0.
\end{equation}
磁场同理满足完全相同的波动方程。与标准形式
$\nabla^2\psi-(1/v^2)\partial^2\psi/\partial t^2=0$ 对照，立即读出
$v=1/\sqrt{\mu_0\varepsilon_0}$。这一推导说明速度结论并不依赖于选取沿 $x$ 传播、沿 $y$ 偏振的
特殊坐标；特殊坐标只是把矢量方程变成容易计算的分量式。

把平面简谐波代回波动方程可得 $\omega^2=c^2k^2$。正频率支对应沿 $+\vect k$ 方向传播，改变相位
符号则得到反向传播的波。波动方程是线性的：若两列波分别是解，它们的和仍是解，因此干涉、驻波和
偏振叠加都自然成立。所谓“振幅任意”只是在给定线性真空方程时，方程本身不替波源规定振幅；实际
系统中的天线功率、边界条件和能量守恒仍会决定可实现的场强。类似地，频率可由波源选择，但波速固定，
于是不同频率对应不同波长。可见光只是电磁波谱的一小段，改变频率并不改变其电磁波本性。

图~\ref{ch26:ref:review-wave-3d} 中电场位于蓝色平面，磁场位于灰色平面；其方向在下文逐项说明。

\planeWaveThreeDDiagram

图中蓝色曲线位于 $x$--$y$ 平面，灰色曲线位于 $x$--$z$ 平面；同一 $x$ 处的两根场矢量同相。
右手从 $\vect E$ 旋向 $\vect B$，拇指所指就是传播方向。图上的斜向 $z$ 轴只是纸面投影，不表示
磁场真的“倾斜”到传播方向。若 $\vect E$ 始终沿固定横向振动，这列波为线偏振波；若两个横向分量
存在相位差，电场矢量端点可描出椭圆或圆。


\subsection{场的能量与坡印廷矢量}


真空中电场与磁场的能量密度分别为
\begin{equation}
  u_E=\frac12\varepsilon_0E^2,
  \qquad
  u_B=\frac{B^2}{2\mu_0},
  \qquad u=u_E+u_B.
\end{equation}
为了找出能量怎样流动，将 $\vect E$ 点乘安培--麦克斯韦方程、将 $\vect B/\mu_0$ 点乘法拉第方程，
除去公共常数后，两个点乘式分别为
\begin{align}
 \frac1{\mu_0}\vect E\cdot(\nabla\times\vect B)
   &=\vect j\cdot\vect E+\frac{\partial}{\partial t}
      \left(\frac{\varepsilon_0E^2}{2}\right),\\
 \frac1{\mu_0}\vect B\cdot(\nabla\times\vect E)
   &=-\frac{\partial}{\partial t}\left(\frac{B^2}{2\mu_0}\right).
\end{align}
用第二式减去第一式，再利用恒等式
\begin{equation}
  \nabla\cdot(\vect E\times\vect B)
  =\vect B\cdot(\nabla\times\vect E)-
  \vect E\cdot(\nabla\times\vect B),
\end{equation}
可得
\begin{equation}
  \boxed{\frac{\partial u}{\partial t}+\nabla\cdot\vect S=-\vect j\cdot\vect E},
  \qquad
  \boxed{\vect S=\frac{1}{\mu_0}\vect E\times\vect B}.
\end{equation}
这就是电磁场的局部能量守恒式。$\vect S$ 称为坡印廷矢量，它的方向是能量传播方向，大小表示
单位时间穿过单位垂直面积的能量。

静电场和稳恒磁场可以把能量储存在空间中；当电场与磁场共同组成传播的电磁波时，能量还能够从一处
流向另一处。太阳并不需要把物质一路输送到地球，太阳辐射携带的能量正是由电磁场穿过空间传来的。
因此坡印廷矢量不是只为配合波动图形而引入的方向箭头，它给出了能量运输的定量描述。对于平面波，
$\vect E\times\vect B$ 指向传播方向，这也反过来帮助判断两场矢量的相对取向。

对有限体积 $V$ 积分，局部守恒式还可以写成
\begin{equation}
  \frac{\dd}{\dd t}\int_Vu\,\dd V
  =-\oint_{\partial V}\vect S\cdot\dd\vect A-
  \int_V\vect j\cdot\vect E\,\dd V.
\end{equation}
左端是体内电磁场能量的增加率；第一个负号项表示由边界流出的电磁能；第二项是场对带电物质做功的
功率。若 $\vect j\cdot\vect E>0$，场把能量交给物质，例如电阻发热；若它为负，外力驱动电荷并向场
输送能量，例如发射天线。$\vect S$ 的单位为 $\si{W/m^2}$，与“单位面积上的功率”这一物理解释一致。

图~\ref{ch26:ref:review-energy-balance} 中边界向外的箭头对应面积分，体内同向的电场、电流对应正的做功项。

\poyntingBalanceDiagram

对上面的真空平面波，$E=cB$ 且 $c^2=1/(\mu_0\varepsilon_0)$，所以
\begin{equation}
  u_E=\frac12\varepsilon_0E^2
  =\frac{B^2}{2\mu_0}=u_B.
\end{equation}
电、磁两部分能量在每一时刻都相等。瞬时能流密度为 $S=\varepsilon_0cE^2$，一个周期内的平均光强为
\begin{equation}
  \boxed{\langle S\rangle=\frac12\varepsilon_0cE_0^2
  =\frac{cB_0^2}{2\mu_0}}.
\end{equation}


% EM teaching enrichment 20260927
\paragraph{积分形式中的三个能量去向}
对固定体积 $V$，坡印廷定理写成
\[
 -\oint_{\partial V}\boldsymbol S\cdot\dd\boldsymbol A
 =\frac{\dd}{\dd t}\int_Vu\dd V+\int_V\boldsymbol J\cdot\boldsymbol E\dd V.
\]
左边是从边界净流入的功率，右边分别为内部场能增加率与场对物质的做功率。若是稳恒欧姆导体，场能不变，流入能量最终转化为焦耳热；若是真空充电电容器间隙，没有传导耗散，流入主要增加场能。
对平行板半径 $R$、间距 $d$ 的拟稳充电模型，边缘附近 $B=\mu_0I/(2\pi R)$，侧面流入功率的大小为
$(EB/\mu_0)(2\pi Rd)=IEd=IU$，与电路功率一致。此处忽略边缘修正，只用来核对能量流向和总量。


\subsection{平面波方向与相位判读}


在无电荷、无电流的真空区域，对法拉第定律再取旋度并使用
$\nabla\times(\nabla\times\vect E)=\nabla(\nabla\cdot\vect E)-\nabla^2\vect E$，得到
\begin{equation}
 \nabla^2\vect E-\mu_0\varepsilon_0\frac{\partial^2\vect E}{\partial t^2}=0.
\end{equation}
磁场满足同样的波动方程，传播速度为
\begin{equation}
 c=\frac1{\sqrt{\mu_0\varepsilon_0}}.
\end{equation}
平面波中 $\vect E$、$\vect B$ 和传播方向 $\uvec k$ 两两垂直，并满足
\begin{equation}
 \vect B=\frac1c\uvec k\times\vect E,
 \qquad B=\frac Ec.
\end{equation}

\begin{center}
\begin{tikzpicture}[x=.85cm,y=.78cm,>=Latex]
 \foreach \y/\lab in {1.65/{E_x/E_0},-1.65/{cB_y/E_0}}{
  \draw[axis] (0,\y)--(8.3,\y) node[right] {$z$};
  \draw[axis] (0,\y-1.25)--(0,\y+1.3);
  \node[anchor=east] at (-.15,\y+.8) {$\lab$};
  \draw[DeepBlue,very thick] plot[domain=0:8,samples=180] (\x,{\y+sin(90*\x)});
 }
 \foreach \x in {1,3,5,7}{\draw[guide] (\x,-2.8)--(\x,2.8);}
 \draw[->,TextGray,thick] (5,3.2)--(7,3.2) node[midway,above] {传播方向 $+\uvec z$};
\end{tikzpicture}
\figcap{同一时刻的电、磁场分量，分两行绘出以比较相位。}\label{ch26:fig:wave-phase}
\end{center}
图~\ref{ch26:fig:wave-phase} 中两行的波峰、波谷和零点位于相同的 $z$，因此两场同相。
纵坐标分别为归一化分量，而非两个空间坐标；实际方向为 $\vect E\parallel\uvec x$、
$\vect B\parallel\uvec y$，由 $\uvec x\times\uvec y=\uvec z$ 确定向前的能流。


% EM teaching enrichment 20260927
\paragraph{从一个电场分量补齐整束平面波}
给定 $\boldsymbol E=E_0\cos(kz-\omega t)\boldsymbol e_x$，等相位条件 $kz-\omega t=\text{常量}$ 表明波沿 $+z$ 传播。于是
\[
 \boldsymbol B=\frac{E_0}{c}\cos(kz-\omega t)\boldsymbol e_y,
 \qquad \omega=ck.
\]
电场与磁场同时反向时，$\boldsymbol E\times\boldsymbol B$ 仍沿 $+z$。若把相位改为 $kz+\omega t$，传播方向反向，磁场的相对方向也必须反向。
只画两条正弦曲线还不足以说明空间关系；应同时给出坐标轴、传播方向以及叉积检验。波形图纵坐标表示场分量，不是粒子在空间中的上下振动轨迹。


\subsection{能量、动量与辐射压}


电磁场能量密度与坡印廷矢量为
\begin{equation}
 u=\frac12\varepsilon_0E^2+\frac{B^2}{2\mu_0},
 \qquad \vect S=\frac1{\mu_0}\vect E\times\vect B.
\end{equation}
$\vect S$ 给出单位面积单位时间通过的能量及其方向。真空简谐平面波的平均光强
\begin{equation}
 I=\langle S\rangle=\frac12c\varepsilon_0E_0^2.
\end{equation}
正入射完全吸收表面的平均辐射压为 $I/c$，正入射理想反射表面为 $2I/c$；斜入射还需取法向动量变化。

\begin{origexamplebox}{激光束的场强和辐射压}
功率 $P$ 均匀分布在截面积 $A$ 的连续激光束中，将光束近似为单色线偏振平面波。求真空中的电场峰值；若光束正入射并被完全吸收，求平均压力。
\end{origexamplebox}
\begin{xesolution}
光强 $I=P/A$。由 $I=c\varepsilon_0E_0^2/2$，
\[
 \boxed{E_0=\sqrt{\frac{2P}{Ac\varepsilon_0}}}.
\]
峰值电场与有效值电场相差 $\sqrt2$，上式中的因子 $2$ 来自对 $\cos^2\omega t$ 取平均。
在短时间 $\Delta t$ 内，到达面积 $A$ 的能量为 $IA\Delta t$，携带向前的动量为 $IA\Delta t/c$。
完全吸收时这部分动量全部交给物体，故 $F=(IA\Delta t/c)/\Delta t=IA/c$，再除以面积得到压力
\[
 \boxed{p_{\rm rad}=\frac{P}{Ac}}.
\]
若改为理想镜面反射，法向动量反向，压力加倍。$P/(Ac)$ 的单位为 $\mathrm{N/m^2}$；
保持面积不变而使功率趋零，场强和压力也都趋零。
\end{xesolution}

\begin{center}
\begin{tikzpicture}[scale=.92,>=Latex]
  \draw[very thick] (5,-1.5)--(5,1.5);
  \foreach \y in {-1,-.5,0,.5,1}{\draw[vector] (.3,\y)--(3.5,\y);}
  \draw[->,TextGray,very thick] (5,0)--(6.5,0) node[right]{$\vect F$};
  \draw[->,DeepBlue,thick] (3.8,.85)--(4.8,.85) node[midway,above]{$\vect p_\gamma$};
  \node[clearlabel,above] at (5,1.65) {吸收面};
  \node[clearlabel] at (1.9,-1.35) {均匀光束，光强 $I$};
\end{tikzpicture}
\figcap{光束被吸收时把传播方向的动量交给物体；压力等于单位面积单位时间的法向动量通量。}\label{ch26:fig:radiation-pressure}
\end{center}
图~\ref{ch26:fig:radiation-pressure} 中力箭头与入射传播方向同向；若改成反射，光子法向动量反向，力增大一倍。

\begin{origexamplebox}{太阳辐射对完全反射薄膜的推力}
面积为 $A$ 的平面薄膜正对光强为 $I$ 的光束，薄膜作理想镜面反射且光束完全覆盖薄膜。求推力；若光束与薄膜法线夹角为 $\theta$，求法向力。
\end{origexamplebox}
\begin{xesolution}
正入射反射使单位能量携带的法向动量从 $+1/c$ 变为 $-1/c$，变化量为 $2/c$，所以
\[
 \boxed{F=\frac{2IA}{c}}.
\]
斜入射时，设 $0\le\theta\le\pi/2$，薄膜在垂直光束方向的投影面积为 $A\cos\theta$，整片薄膜截获的功率为 $IA\cos\theta$，每单位能量的法向动量变化为
$2\cos\theta/c$，故
\[
 \boxed{F_n=\frac{2IA}{c}\cos^2\theta}.
\]
$\theta=0$ 回到正入射结果，$\theta\to\pi/2$ 时截获功率和法向动量都趋零。$IA/c$ 的单位为牛顿。
\end{xesolution}


% EM structure and teaching revision 20260928
\subsection{场幅、平均强度与总功率的区别}
真空平面波中 $B_m=E_m/c$，电场与磁场能量密度相等，瞬时能流方向为 $\boldsymbol E\times\boldsymbol B$。对于单频正弦波，时间平均强度为
\[
 \langle S\rangle=\frac12\varepsilon_0cE_m^2
 =\varepsilon_0cE_{\mathrm{rms}}^2.
\]
峰值与有效值混用会造成二倍误差。总功率还需对横截面积分：$P=\int\langle\boldsymbol S\rangle\cdot\dd\boldsymbol A$。只有波束在所选截面上近似均匀并正入射时，才可简写 $P=\langle S\rangle A$。
设均匀波束以角 $\theta$ 入射平面，强度按垂直波束方向的单位面积定义，则表面每单位面积截获的能量通量为 $\langle S\rangle\cos\theta$。若镜面理想反射，法向动量反转再提供一个 $\cos\theta$，因而法向辐射压为
$p_n=2\langle S\rangle\cos^2\theta/c$。
斜入射问题若同时漏掉投影面积和动量投影，结果可能在正入射时正确，却在掠入射极限仍给出有限法向压力。用 $\theta\to\pi/2$ 作检查，能直接识别这类错误。


\section{辐射、散射与介质传播}

\subsection{加速电荷的辐射}


静止电荷产生静电场；匀速运动的电荷在固定观察点也可产生随时间变化的场，但在真空中不产生向远区输送净能量的辐射项。电磁作用始终以有限速度传播，不能把“有时间变化”和“有辐射”等同。加速度使场出现按 $1/r$ 衰减的辐射项。在非相对论条件 $v\ll c$ 下，距离电荷很远处的辐射电场大小为
\begin{equation}
  \boxed{E_{\rm rad}(r,t)=\frac{1}{4\pi\varepsilon_0}
  \frac{|q|\,a(t-r/c)\sin\theta}{c^2r}}.
\end{equation}
其中 $t-r/c$ 是推迟时刻，$\theta$ 是观察方向与加速度方向的夹角。这里用 $E_{\rm rad}$ 而不用
$E_r$：辐射场实际上垂直于传播方向，不是沿半径方向的场分量。
$a$ 在大小公式中表示加速度的模；方向应由矢量式判断：
\[
 \vect E_{\rm rad}
 =\frac{q}{4\pi\varepsilon_0c^2r}\,
 \uvec n\times(\uvec n\times\vect a)
 \quad\text{（取推迟时刻）}.
\]
利用三重积，$\uvec n\times(\uvec n\times\vect a)=
\uvec n(\uvec n\cdot\vect a)-\vect a$，恰好是加速度横向分量的反向；
电荷为负时场方向再反转。

辐射场随 $1/r$ 衰减，而库仑近场随 $1/r^2$ 衰减；所以在足够远处，辐射项反而占主导。远区有
$B_{\rm rad}=E_{\rm rad}/c$，于是单位立体角内的辐射功率为
\begin{equation}
  \frac{\dd P}{\dd\Omega}=r^2S
  =\frac{q^2a^2\sin^2\theta}{16\pi^2\varepsilon_0c^3}.
\end{equation}

图~\ref{ch26:ref:review-radiation-pattern} 中角度 $\theta$ 从加速度轴量起，蓝色曲线的径向长度代表功率而不是电场幅值。
图中另画出的 $\vect E_{\rm rad}$ 箭头对应 $q>0$：它垂直于观察方向 $\uvec n$，
沿 $\uvec n\times(\uvec n\times\vect a)$ 指向右下方；若电荷为负，电场方向反转而功率角分布不变。

\radiationPatternDiagram

沿加速度方向 $\theta=0,\pi$ 时不辐射，垂直于加速度时最强。对全空间积分，利用
\begin{equation}
  \int\sin^2\theta\,\dd\Omega
  =2\pi\int_0^\pi\sin^3\theta\,\dd\theta=\frac{8\pi}{3},
\end{equation}
得到拉莫尔公式
\begin{equation}
  \boxed{P=\frac{q^2a^2}{6\pi\varepsilon_0c^3}}.
\end{equation}
它说明辐射功率正比于加速度平方；这里的公式只适用于非相对论点电荷，不能直接套用于
$v$ 接近 $c$ 的情形。

图~\ref{ch26:fig:s41-10-7} 用一组相隔半波长的波前表示振荡电荷向外传递的场。每个蓝箭头都沿波前切向，因而与能流
方向垂直；相邻半波带的箭头反向，对应简谐场的相位翻转。沿加速度轴箭头趋于零，与
$E_{\rm rad}\propto\sin\theta$ 一致。

\radiationWavefrontDiagram

辐射图样常画成绕加速度轴的“甜甜圈”。图中的径向长度表示该方向的相对辐射强弱，并不是电荷真的
喷出一圈物质。把 $E_{\rm rad}$ 代入 $S=\varepsilon_0cE_{\rm rad}^2$，其中 $E_{\rm rad}^2\propto1/r^2$；
再乘球面面积元 $r^2\dd\Omega$，距离因子正好消去。因此无论在多大的球面上统计，只要球面包围波源，
在对应同一推迟时刻的波前上，穿出的总辐射功率都相同。若加速度随时间变化，不同半径处同一观察时刻所收到的是不同发射时刻的能量，不能直接要求瞬时功率相等。这是能量从波源持续向外传播而不是在空间中凭空损失的表现。

推迟时刻 $t-r/c$ 是有限传播速度的直接结果：观察者接收到的不是电荷“此刻”的运动信息，而是场从
电荷传播到观察点之前的状态，正如地球上看到的太阳是约八分钟以前的太阳。拉莫尔公式曾是经典
电磁学的重大成功，却也暴露出经典原子模型的根本困难：若电子像行星一样绕原子核转动，它始终有
向心加速度，便应不断辐射并损失机械能，轨道逐渐缩小直至落入原子核。真实原子却稳定存在，这个冲突
不能靠经典电磁学自身解决，后来成为量子理论产生的重要背景之一。

还要分清近场与辐射场。靠近电荷时，$1/r^2$ 的库仑型项以及与速度有关的感应项可能占主导，
它们储存的能量可在波源附近往返交换；远离波源后，按 $1/r$ 衰减的辐射项成为主导，并携带净能量
离开系统。因此“有时变电场”不等于“观察点处一定只有辐射场”，判断远区辐射必须同时看距离依赖和
能流方向。拉莫尔公式计算的是已经离开波源、穿过远处闭合面的总功率。


% EM teaching enrichment 20260927
\paragraph{近场与辐射场的距离检验}
库仑型近场随 $1/r^2$ 衰减，非相对论辐射场随 $1/r$ 衰减。远区能流与场强平方成正比，球面面积又正比于 $r^2$，因此只有 $1/r$ 的场项能在远处留下有限的总辐射功率。
沿加速度轴方向，辐射场的横向投影为零；垂直轴方向最强。图中的 $\sin^2\theta$ 表示每单位立体角的功率分布，不是电场振幅，振幅对应 $\sin\theta$。
若同一非相对论点电荷的加速度大小加倍，拉莫尔功率变为四倍；恒速直线运动的加速度为零，不因观察点处近场随时间改变就产生同样的辐射项。


\subsection{汤姆孙散射截面}


考虑沿 $+x$ 传播、沿 $z$ 偏振的电磁波，对初始平均静止的自由电子作弱场、非相对论近似
\begin{equation}
  E_z=E_0\cos(kx-\omega t).
\end{equation}
电子的振动速度幅值约为 $v_0=eE_0/(m\omega)$。要求 $eE_0/(m\omega c)\ll1$，既保证 $v_0\ll c$，也保证位移幅值 $v_0/\omega$ 远小于 $c/\omega$，可把电子所在区域的电场看成均匀。忽略磁力相对于电力的
$v/c$ 阶小量，电子加速度为
\begin{equation}
  a_z=-\frac{eE_0}{m}\cos\omega t.
\end{equation}
把它代入拉莫尔公式并对周期取平均，因 $\langle\cos^2\omega t\rangle=1/2$，有
\begin{equation}
  \langle P_e\rangle=\frac{e^4E_0^2}{12\pi\varepsilon_0m^2c^3}.
\end{equation}
入射波的平均能流密度是 $\langle S_{\rm in}\rangle=\varepsilon_0cE_0^2/2$。散射截面定义为
“总散射功率除以入射能流密度”，于是
\begin{align}
  \sigma_T
  &=\frac{\langle P_e\rangle}{\langle S_{\rm in}\rangle}
  =\frac{e^4}{6\pi\varepsilon_0^2m^2c^4}
  \\
  &=\boxed{\frac{8\pi}{3}\left(\frac{1}{4\pi\varepsilon_0}
  \frac{e^2}{mc^2}\right)^2}
  =\frac{8\pi}{3}r_e^2
  \approx\SI{6.65e-29}{m^2},
\end{align}
其中 $r_e=e^2/(4\pi\varepsilon_0mc^2)$ 称为经典电子半径。截面具有面积量纲，但它不意味着电子是
半径 $\sqrt{\sigma_T/\pi}$ 的实心小球；它衡量的是电子从入射波中截获并再辐射能量的能力。

在这一低频、非相对论经典近似中，$E_0^2$ 在散射功率与入射能流之比中消去，结果也不显含频率，
所以汤姆孙截面与入射波强度、频率均无关。这个结论有明确的适用边界：频率高到必须考虑电子
反冲和量子效应时（光子能量不再满足 $\hbar\omega\ll mc^2$），需要使用更一般的散射理论，不能仍把自由电子当成被均匀电场驱动的经典点电荷。这里还假定电子不受原子束缚；束缚电子在低频、近共振等区域一般不满足这一频率无关结论。
图中的电子电荷为 $-e$，所以瞬时加速度与电场反向；散射强弱取决于 $a^2$，不因这种反向而改变。

散射截面的定义可以想成：若入射能流密度为 $S_{\rm in}$，把总散射功率写成
$P_{\rm sc}=S_{\rm in}\sigma$，那么 $\sigma$ 就是一个等效受光面积。它允许我们不追踪每一束散射波，
而用一个具有面积量纲的量比较相互作用强弱。对一个自由电子，偶极辐射在电子加速度方向上最弱、在
垂直方向上最强；把各方向功率全部积分才得到上面的 $\sigma_T$。经典电子半径 $r_e$ 也不是测得的
电子几何边界，而是由 $e,m,c,\varepsilon_0$ 组合出来的电磁作用长度尺度。

图~\ref{ch26:ref:review-thomson} 中向下的加速度与向上的瞬时电场反向；两者沿同一振动轴，散射功率取决于其横向投影。

\thomsonScatteringDiagram


% EM teaching enrichment 20260927
\paragraph{截面的面积单位与物理含义}
散射截面定义为总散射功率除以入射光强，因此单位是面积，但它不是把电子当作半径固定的小硬球得到的几何面积。
经典弱场模型中，电子加速度振幅 $a_0=eE_0/m_e$，平均辐射功率正比于 $a_0^2$，入射光强正比于 $E_0^2$，相除后振幅抵消，得到与入射强度无关的截面。
若频率接近束缚电子的共振，或光子反冲与电子静能不可忽略，自由、非相对论电子的这套推导便不足以使用；它不能无条件推广为所有物质、所有频率的散射规律。


\subsection{介质中的传播与色散}


均匀、线性、各向同性且无耗的介质中 $\vect D=\varepsilon\vect E$、$\vect B=\mu\vect H$，相速度
$v=1/\sqrt{\varepsilon\mu}$。实际介质的 $\varepsilon$ 随频率变化，因此折射率
$n(\omega)=c/v_p$ 具有色散。波包包络以群速度
\begin{equation}
 v_g=\frac{\dd\omega}{\dd k}
\end{equation}
传播；在强色散和吸收区，群速度不能不加分析地解释为信息或能量速度。

法向入射到两个无耗、非磁介质的平面波，在界面处切向 $E,H$ 连续。定义波阻抗
\begin{equation}
 Z=\sqrt{\frac{\mu}{\varepsilon}},
\end{equation}
电场振幅反射系数和透射系数为
\begin{equation}
 r_E=\frac{Z_2-Z_1}{Z_2+Z_1},
 \qquad
 t_E=\frac{2Z_2}{Z_1+Z_2}.
\end{equation}
这些振幅系数来自两个边界方程。统一以切向电场的同一正方向记有符号振幅，
则 $E_i+E_r=E_t$；由于反射波传播反向，其磁场对该正方向的贡献反号，故
\[
 \frac{E_i-E_r}{Z_1}=\frac{E_t}{Z_2}.
\]
代入 $E_t=E_i+E_r$，整理为 $(Z_2-Z_1)E_i=(Z_1+Z_2)E_r$，
得到 $r_E$，再用 $t_E=1+r_E$ 得到透射系数。
平均能流为 $E_0^2/(2Z)$，所以
\[
 R=|r_E|^2,\qquad T=\frac{Z_1}{Z_2}|t_E|^2
   =\frac{4Z_1Z_2}{(Z_1+Z_2)^2},\qquad R+T=1.
\]
最后一步是将分子 $(Z_2-Z_1)^2+4Z_1Z_2$ 合并为 $(Z_1+Z_2)^2$。

\begin{center}
\begin{tikzpicture}[scale=.9,>=Latex]
  \fill[DeepBlue!5] (4,-2)--(8,-2)--(8,2)--(4,2)--cycle;
  \draw[very thick] (4,-2)--(4,2);
  \draw[vector] (.5,.6)--(3.65,.6) node[midway,above]{$\vect k_i$};
  \draw[vector] (3.55,-.5)--(.7,-.5) node[midway,below]{$\vect k_r$};
  \draw[vector] (4.35,.6)--(7.5,.6) node[midway,above]{$\vect k_t$};
  \node[clearlabel] at (2,1.55) {$Z_1$};
  \node[clearlabel] at (6,1.55) {$Z_2$};
  \node[clearlabel,below] at (4,-2) {界面};
\end{tikzpicture}
\figcap{法向入射的反射与透射。箭头表示传播方向；反射电场的正负由振幅系数 $r_E$ 决定。}\label{ch26:fig:interface}
\end{center}
图~\ref{ch26:fig:interface} 把三个传播方向分开；$r_E<0$ 表示电场相位翻转，并不把反射箭头改画成向右。


% EM teaching enrichment 20260927
\paragraph{相速度与群速度怎样从同一色散关系得到}
在窄带、弱吸收条件下写 $k=n(\omega)\omega/c$，则
\[
 v_p=\frac c{n(\omega)},\qquad
 v_g=\left(\frac{\dd k}{\dd\omega}\right)^{-1}
 =\frac c{n+\omega\,\dd n/\dd\omega}.
\]
无色散时两者相同；有色散时，单色波峰的速度与包络速度不同。若在某一频段近似 $n=n_0+a\omega$，则 $v_g=c/(n_0+2a\omega)$，并非简单的 $c/n$。
这个线性近似只在所给窄频段使用，不能外推到任意频率。强吸收引起波包重塑时，包络峰值运动也不能直接当作新信息的传播速度。

\begin{origexamplebox}{介质界面的能流守恒}
真空中的平面波法向入射到折射率为 $n$、磁导率近似为 $\mu_0$ 的无耗介质。设 $n>0$。求电场振幅反射系数、功率反射率与透射率，验证能流守恒，并检验 $n=1$ 与 $n\to\infty$ 的极限。
\end{origexamplebox}
\begin{xesolution}
界面没有自由面电流时，切向电场和磁场连续。把入射电场振幅归一化为1，
反射、透射振幅分别为 $r_E,t_E$，则
\[
 1+r_E=t_E,\qquad \frac{1-r_E}{Z_0}=\frac{t_E}{Z_0/n}.
\]
第二式中的负号来自反射波传播方向反转，不能把反射磁场仍按入射方向相加。
消去 $t_E$，即 $1-r_E=n(1+r_E)$。非磁介质波阻抗 $Z=Z_0/n$，于是
\[
 r_E=\frac{Z_2-Z_1}{Z_2+Z_1}
 =\frac{1-n}{1+n}.
\]
因此
\[
 \boxed{R=|r_E|^2=\left(\frac{n-1}{n+1}\right)^2}.
\]
透射振幅系数 $t_E=1+r_E=2/(1+n)$，所以
\[
 T=n|t_E|^2=\frac{4n}{(1+n)^2},\qquad
 R+T=\frac{(n-1)^2+4n}{(n+1)^2}=1.
\]
只有 $n>1$ 时 $r_E<0$，反射电场相位翻转；$0<n<1$ 时不翻转。
$n=1$ 时 $R=0,T=1$；在本无耗实折射率模型的 $n\to\infty$ 极限中 $R\to1,T\to0$。
后一极限不等于对真实吸收金属的描述。$R,T$ 均无量纲且位于 $[0,1]$。
\end{xesolution}


% EM structure and teaching revision 20260928
\subsection{平面波界面的振幅与能流核对}
两种无耗散、非磁性介质折射率为 $n_1,n_2$，平面波正入射，取同一切向电场正向。电、磁场的切向边界条件给出
\[
 E_i+E_r=E_t,\qquad n_1(E_i-E_r)=n_2E_t.
\]
反射波传播方向反转，故磁场方程中的反射项带负号。解得电场振幅系数
$r=(n_1-n_2)/(n_1+n_2)$、$t=2n_1/(n_1+n_2)$。负的 $r$ 表示电场相位反转，不表示能量反向增加。
能流与 $nE^2$ 成正比，所以反射率与透射率为
\[
 \mathcal R=|r|^2,\qquad \mathcal T=\frac{n_2}{n_1}|t|^2,
 \qquad \mathcal R+\mathcal T=1.
\]
例如空气近似 $n_1=1$、另一介质 $n_2=1.5$，有 $r=-0.2$、$t=0.8$，反射能量占 $4\%$，透射占 $96\%$。把 $t^2=0.64$ 当成透射率会误判能量不守恒。
这些公式使用正入射、线性、无耗散且磁导率相同的前提。斜入射须区分偏振；存在吸收时还要计入介质内的能量沉积。


\section{超导磁响应与磁通量子化}

\subsection{超导：零电阻还不够}


金属电阻率一般随温度降低而下降，因为晶格热振动减弱，传导电子受到的散射减少。随着低温技术发展，
人们自然追问电阻会趋向一个非零常数，还是能够真正消失。1908 年，荷兰物理学家卡末林·昂内斯（图~\ref{ch26:ref:review-onnes-portrait}）
成功液化氦，获得约 $\SI{4.2}{K}$ 的低温条件；1911 年他研究汞的电阻时，发现温度降到这一范围后，
电阻突然跌至实验分辨率以下，超导态由此被发现。

\historyfigure[0.23\linewidth]{onnes-built-in-v1.png}%
{卡末林·昂内斯（Heike Kamerlingh Onnes），低温物理和超导现象的先驱；历史人物肖像。\label{ch26:ref:review-onnes-portrait}}\label{ch26:fig:s41-10-8}

后来人们在多种元素、合金和铜氧化物中发现超导转变。材料进入超导态的温度称为临界温度 $T_c$；
实际超导还受临界磁场和临界电流密度限制，不能只凭 $T_c$ 判断应用性能。从 1911 年的汞到 1973 年前后，
传统材料的 $T_c$ 提升十分缓慢；1986 年贝德诺尔茨（Johannes Georg Bednorz）与缪勒（Karl Alexander M\"uller）发现铜氧化物高温超导体，随后材料的
临界温度迅速越过液氮沸点 $\SI{77}{K}$。下图以分材料族的历史数据展示不同年代的临界温度，反映了
超导材料探索中长期缓慢推进与关键发现后快速跃升并存的过程。

图~\ref{ch26:ref:review-superconductivity-timeline} 串起零电阻的发现、迈斯纳效应和新材料温区，图中的年份用于定位各项发现。

\superconductivityTimeline

\superconductivityChartFaithful

1987 年，吴茂昆、朱经武等人的钇钡铜氧材料使转变温度达到约 $90\,$K，超过液氮沸点。随后铋系、铊系和汞系铜氧化物又推动了高温区的探索。早期有过常规超导临界温度约 $30\,$K 的经验估计，但它不是 BCS 理论给出的普适严格上限，不能把后来的高温发现理解为违反一条已证明的上限定理。

零电阻最直接的后果，是闭合超导环中的电流可以长期存在而不发生焦耳耗散。但若只把超导体看作
电阻率严格为零的普通导体，仍解释不了它最重要的磁学性质。

“电阻测不到”与“电阻严格为零”在实验上并不完全等价，因为任何电压表都有分辨率。超导持续电流
实验提供了另一种证据：在闭合超导环中建立电流后，若电阻哪怕很小，磁能也应按一定时间尺度衰减；
实际观测到电流可以在极长时间内保持，由衰减上限反推的电阻率远低于普通良导体。零电阻意味着
$P_J=I^2R$ 不再产生焦耳热，却不意味着超导装置可以承受任意电流和磁场。温度、外磁场或电流密度
超过相应临界条件时，材料会回到有耗散的正常态。

超导临界温度的发展也说明实验条件与材料体系同样重要。早期研究依赖液氦低温，能探索的材料范围和
温区都很有限；传统金属与合金的转变温度长期只缓慢提高。铜氧化物高温超导出现后，液氮温区的实验
和应用成为可能，材料的晶体结构、载流子浓度及制备工艺也成为研究重点。图~\ref{ch26:fig:s41-10-9} 将材料名称、年代和实验条件列在各行，用横向长度表示临界温度，避免低温点拥挤。读图时应比较同一时代的数量级与跃迁趋势，
而不宜只盯住某一个最高点；图中高压条件下的数据也不能与常压工程材料直接等同。


% EM teaching enrichment 20260927
\paragraph{持续电流能给出多小的电阻上限}
把环近似为自感 $L$ 与极小电阻 $R$ 的闭合回路，则 $I(t)/I_0=e^{-Rt/L}$。若观察时间 $t$ 内电流相对下降不超过 $\delta$，便有
\[
 R\le-\frac Lt\ln(1-\delta)\simeq\frac{L\delta}{t}\quad(\delta\ll1).
\]
这是一条由测量精度得到的上限，不是由有限精度直接证明一个数学量严格为零。它说明长时间稳定的持续电流为何是灵敏的实验检验，同时也提示应把“测不到变化”与“没有任何变化”在表述上分开。


\subsection{理想导体与迈斯纳效应}


在理想导体内部，$\vect E=0$。法拉第定律给出
\begin{equation}
  \frac{\dd\Phi_B}{\dd t}=-\oint_C\vect E\cdot\dd\vect l=0.
\end{equation}
因此理想导体只保证磁通不再改变：降温前内部若已有磁场，进入零电阻状态后可以把这部分磁通“冻结”
下来。超导体却表现出更强的性质——在给定温度、外场条件下进入平衡迈斯纳态时，体内磁场被排斥到表面附近，这就是迈斯纳效应。对第二类超导体，需要外场低于相应下临界场；处于混合态时有量子化涡旋进入，不能笼统断言所有超导状态的体内场都为零。

两种操作顺序最能看出差别。若先把无磁场中的样品降至零电阻态、再施加外场，理想导体中的感应电流
会抵消磁通变化，体内可以保持 $B=0$；这个结果尚不能区分理想导体和超导体。若先在正常态样品中
施加磁场，再降温穿过 $T_c$，理想导体模型预言原有磁通应被冻结在体内。1933 年迈斯纳与奥森菲尔德
测量锡等样品的磁场分布，却发现无论先加磁场还是先降温，只要样品进入超导态且外场低于临界条件，
在平衡迈斯纳态下体内磁场都会被排斥。因而迈斯纳效应是一种平衡态磁学性质，不只是“电阻为零导致磁通不变”。

从磁场边界条件看，体内近似 $B=0$ 必须由表面附近的屏蔽电流来维持。这个电流不是普通电阻回路中
由持续电场推动的电流，而是超导凝聚态对外磁场的无耗散响应。宏观图常把磁感线画成在几何表面突然
转弯，真实情况则是在伦敦穿透深度的尺度上连续衰减。样品尺寸远大于穿透深度时，把体内看成
$B=0$ 是很好的竞赛近似；研究薄膜或表面效应时，就必须保留这段有限厚度。

图~\ref{ch26:ref:review-meissner} 中左图保留先前的内部磁场，右图外部场线绕开样品；蓝色场线彼此不相交，并满足样品表面无磁通穿入的条件。

\meissnerComparisonDiagram

图中左侧的磁力线穿过材料，表示理想导体保留初始磁通；右侧磁力线绕开超导体，表示迈斯纳排磁。
真实超导体的磁场并非在数学意义上于表面瞬间降为零，而是在称为伦敦穿透深度的薄层内衰减；
在本章的宏观近似中，可写成超导体体内 $B\approx0$。

因此，“完美导电性”和“完全抗磁性”回答的是两个不同问题。前者回答电场能否在体内维持以及电流是否
耗散；后者回答给定温度和外磁场后，平衡态磁通怎样分布。理想导体的最终磁场依赖进入零电阻状态前的
历史，超导体的迈斯纳态则在相同外部条件下趋向同一排磁结果。正是这种与路径无关的排磁性，使超导体
不能被简单等同为电导率无限大的普通金属。

把磁体放在超导体上方时，表面屏蔽电流改变空间磁场分布，可产生向上的磁力。若磁力与重力平衡，
磁体能够悬浮；图~\ref{ch26:ref:review-levitation-scene} 中磁体与低温样品间的间隙体现这一状态。实验中常见的稳定悬浮还与第二类超导体的磁通钉扎有关。

图~\ref{ch26:ref:review-levitation} 的右侧是磁体的受力图：向上的磁力与向下的重力都作用于该磁体，
平衡时二者大小相等。下方超导圆盘受到的力不属于这个受力分析对象。

\levitationDiagram

\historyfigure[0.32\linewidth]{levitation-built-in-v1.png}%
{超导磁悬浮实验：磁场与屏蔽电流共同产生向上的电磁力。\label{ch26:ref:review-levitation-scene}}


\subsection{超导环中的磁通保持}


以下采用宏观厚环、忽略动能电感的近似，并要求全过程不失超、无相位滑移且量子数保持不变。设超导环的几何自感为 $L$，外界穿过环的磁通为 $\Phi_{\rm ext}$，环中持续电流 $I$ 自身贡献磁链 $LI$。
环路电阻为零时，感应电动势满足
\begin{equation}
  L\frac{\dd I}{\dd t}=-\frac{\dd\Phi_{\rm ext}}{\dd t},
\end{equation}
所以
\begin{equation}
  \boxed{\frac{\dd}{\dd t}\bigl(\Phi_{\rm ext}+LI\bigr)=0}.
\end{equation}
先在外磁场中把圆环冷却到超导态，再撤去外场，$\Phi_{\rm ext}$ 减小；环中便出现持续电流，使 $LI$
增加，从而补偿总磁通的变化。

若圆环建立该超导状态时的总磁通为 $\Phi_{\rm ini}$，以后在上述适用条件内缓慢改变外磁场，都有
\begin{equation}
  I=\frac{\Phi_{\rm ini}-\Phi_{\rm ext}}{L}.
\end{equation}
外磁通减小时电流按楞次定律方向增大，外磁通增大时电流反向调整。这里保持的是“外磁通加自感磁通”
的总和，不是说外磁场本身被冻结不变。对于尺度足够大的经典问题，这一磁通守恒关系直接来自零电阻
环路的法拉第定律；在微观超导理论中还会出现磁通量子化，那是比本章宏观模型更进一步的内容。

图~\ref{ch26:ref:review-ring-hold} 中外场与自感场均取出纸面为正；逆时针持续电流按右手定则补偿逐渐撤去的外磁通。

\superconductingRingDiagram

这个宏观推导抓住了固定状态分支中总磁链保持的物理图像。$\Phi_{\rm ini}$ 是初始总磁通，不是后文定义的磁通量子 $\Phi_0=h/(2e)$。后面的磁通量子化一节进一步区分总磁通和磁通子，并讨论固定状态与最低能量状态的区别。

% EM teaching enrichment 20260927
\paragraph{撤场以后能量从何而来}
在固定状态分支中，若起初环流为零、外磁通为 $\Phi_i$，随后外磁通降到零，则
$I_f=\Phi_i/L$，最终环内能量为 $\Phi_i^2/(2L)$。
这不是零电阻环无中生有地产生能量；撤场装置与环之间通过互感交换功，必须把外磁场源及其驱动过程一并计入能量账。
零电阻保证没有 $I^2R$ 耗散，却不保证环的磁能不变。若改变几何结构而使 $L$ 改变，也应对 $\Phi_{\rm ext}+LI$ 整体求导，不能继续把 $L$ 提到导数外。



\subsection{超导环的量子态选择}


超导体进入迈斯纳态时，体内磁感应强度被排斥；含孔的超导环还允许磁通以离散方式被俘获。对电荷为 $2e$ 的库珀对，磁通量子为
\begin{equation}
 \Phi_0=\frac{h}{2e}.
\end{equation}
在壁厚远大于穿透深度、环路位于超导体内部的理想情形，磁通子（fluxoid）量子化可近似写成
$\Phi=n\Phi_0$。薄环或沿积分路径的超电流密度不可忽略时，量子化的是
\begin{equation}
 \Phi+\mu_0\lambda_L^2\oint_C\vect j_s\cdot\dd\vect l=n\Phi_0.
\end{equation}
其中 $\lambda_L$ 为伦敦穿透深度，$\vect j_s$ 为超电流密度，$C$ 位于超导材料内，磁通 $\Phi$ 的正法向与 $C$ 按右手规则配对。
第二项反映载流子运动；只有选取深入厚壁、$\vect j_s\approx0$ 的路径时才能略去。
不能把任意几何区域的磁通都机械地取整。用集中参数描述薄环时，可将这一项写成 $L_kI$，
并用总电感 $L=L_g+L_k$ 代替单纯的几何自感。

\begin{center}
\begin{tikzpicture}[scale=.9,>=Latex]
 \fill[DeepBlue!12,even odd rule] (0,0) circle (1.7) (0,0) circle (.9);
 \draw[thick] (0,0) circle (1.7) (0,0) circle (.9);
 \foreach \x/\y in {0/0,.45/0,-.45/0,0/.45,0/-.45}{
   \draw[DeepBlue] (\x,\y) circle (.09);\fill[DeepBlue] (\x,\y) circle (.025);
 }
 \draw[->,DeepBlue,very thick] (25:1.3) arc[start angle=25,end angle=150,radius=1.3];
 \node[above] at (0,1.9) {$I>0$：逆时针};
 \node[align=left,anchor=west] at (2.4,.3)
  {俯视图：$\odot$ 表示出纸面磁通\\[5pt]
   取出纸面为正，$I>0$ 产生正自感磁通\\[5pt]
   $I<0$ 时环流和自感场同时反向};
\end{tikzpicture}
\figcap{含孔超导环中的俘获磁通。环流方向调节总磁通子满足量子化条件。}\label{ch26:fig:flux-quantization}
\end{center}
图~\ref{ch26:fig:flux-quantization} 中持续电流产生的自感磁通与外磁通相加，正是下题模型 $\Phi_{\rm ext}+LI=n\Phi_0$ 的几何含义。

\begin{origexamplebox}{超导环的最邻近磁通态}
超导环外加磁通为 $\Phi_{\rm ext}$，总电感为 $L$（厚环可忽略动能电感）。假设不同量子态之间能够充分弛豫，并忽略其他态间能量差。用简化模型
$\Phi_{\rm ext}+LI=n\Phi_0$，相应电感能量为 $U_n=LI^2/2$。求给定 $\Phi_{\rm ext}$ 时的平衡量子数选择规则，并求相邻两态简并的位置。
\end{origexamplebox}
\begin{xesolution}
第 $n$ 态的持续电流和能量为
\[
 I_n=\frac{n\Phi_0-\Phi_{\rm ext}}{L},
 \qquad
 U_n=\frac{(n\Phi_0-\Phi_{\rm ext})^2}{2L}.
\]
平衡态选择使 $|n\Phi_0-\Phi_{\rm ext}|$ 最小的整数 $n$，即
\[
 \boxed{n=\operatorname*{argmin}_{m\in\mathbb Z}|m-\Phi_{\rm ext}/\Phi_0|}.
\]
$n$ 与 $n+1$ 两态简并时
\[
 (n\Phi_0-\Phi_{\rm ext})^2=[(n+1)\Phi_0-\Phi_{\rm ext}]^2,
\]
两边相减，得
\[
 [(n+1)-n]\Phi_0[(2n+1)\Phi_0-2\Phi_{\rm ext}]=0.
\]
因 $\Phi_0\ne0$，故
\[
 \boxed{\Phi_{\rm ext}=\left(n+\frac12\right)\Phi_0}.
\]
$\Phi/L$ 的单位为安培，能量式单位为焦耳。远离半整数点时最低能态唯一；在简并点两种环流分别为 $-\Phi_0/(2L)$ 和 $+\Phi_0/(2L)$，能量均为 $\Phi_0^2/(8L)$。
上述规则选择全局最低能量态，并不保证实际圆环随外场改变就立即跳到该态。若量子态转换需要跨越相位滑移势垒，圆环可滞留在固定 $n$ 的亚稳态；此时应沿同一个 $I_n(\Phi_{\rm ext})$ 分支计算，直到发生允许的状态转换。
\end{xesolution}

\begin{starread}{波导、驻波与边界条件}
理想导体表面切向电场为零。空心、单连通截面的金属波导只允许满足边界条件的离散横向模式，并存在截止频率。
这里不包括同轴线等双导体结构；后者可以支持无截止频率的横电磁（TEM）模式。
对真空填充、理想无耗的矩形波导，横电（TE）与横磁（TM）模式具有不同场结构；传播模式满足 $\omega^2=c^2k_z^2+\omega_c^2$，故 $v_p=\omega/k_z>c$、$v_g=c^2k_z/\omega<c$，并不允许超光速传信。
谐振腔把传播方向也离散化，品质因数由单位周期储能与损耗之比决定。电磁模式量子化后，每个模式成为频率为
$\omega$ 的量子谐振子，这是黑体辐射与光子概念的基础之一。
\end{starread}

% EM structure and teaching revision 20260928
\subsection{磁通保持、磁通子与材料临界条件}
理想导体由零电阻导出的结论是既有磁通不随时间改变；超导体的平衡磁响应还包括迈斯纳效应。二者对“先加磁场后降温”和“先降温后加磁场”的响应不能简单等同。
闭合超导环的量子条件约束的是磁通子。若把几何电感与动能电感之和记为 $L_{\mathrm{tot}}$，外磁通为 $\Phi_{\mathrm{ext}}$，简化模型给出
\[
 \Phi_{\mathrm{ext}}+L_{\mathrm{tot}}I=n\Phi_0,\qquad
 W_n=\frac{(n\Phi_0-\Phi_{\mathrm{ext}})^2}{2L_{\mathrm{tot}}},\quad
 \Phi_0=\frac{h}{2e}.
\]
相邻两条能量抛物线在半整数外磁通处相交，比较 $W_n$ 可确定热力学最低能态。但实际扫场时，能垒会使环停留在亚稳态，未必立即跳到最低的一支。
由某个整数 $n$ 算得电流之后，还须核对临界电流、临界磁场和温度条件；若超出材料允许范围，零电阻与给定电感的模型就不能继续使用。把“允许的量子态”“最低能态”与“实验能到达的状态”分开，才不会把整数取整规则误作所有情形下的实际演化规律。
}

\endgroup
